\documentclass[11pt]{article}
\usepackage[margin=1in]{geometry}
\usepackage[T1]{fontenc}
\usepackage[utf8]{inputenc}
\usepackage[french]{babel}
\usepackage{amsmath,amssymb,amsthm}
\usepackage[colorlinks=true,linkcolor=blue,urlcolor=blue]{hyperref}
\newtheorem{problem}{Problème}
\newenvironment{refs}{\par\small\noindent\emph{Références :}\begin{itemize}\setlength\itemsep{0pt}}{\end{itemize}\normalsize}
\title{Hors de la Caverne \\ \Large Topologie}
\author{}
\date{}
\begin{document}
\maketitle
\noindent\emph{Des problèmes, pas de réponses.} Chaque problème ci-dessous est posé
sans solution. Les références indiquent où trouver les connaissances nécessaires.
\medskip


\section*{Lycée}

\begin{problem}[{Les ponts de Königsberg --- Lycée}]
\textbf{TOP-01.} La ville de Königsberg comptait sept ponts reliant deux îles fluviales aux deux
rives de la Pregel : la plus grande île était reliée par deux ponts à chaque rive et par un pont à la
plus petite île, elle-même reliée par un pont à chaque rive. Les habitants se demandaient si l'on
pouvait faire une promenade traversant chaque pont exactement une fois.
\begin{enumerate}
\item[(a)] Démontrez qu'aucune telle promenade n'existe.
\item[(b)] Modélisez la situation par un graphe --- les masses de terre comme sommets, les ponts comme arêtes
--- et déterminez, avec démonstration, exactement quels graphes connexes admettent une promenade
empruntant chaque arête une fois : en fonction du nombre de sommets rencontrant un nombre impair
d'arêtes, quand une telle promenade existe-t-elle, et quand peut-elle en outre revenir à son point de
départ ?
\end{enumerate}

\end{problem}

\noindent Euler a réglé la question en 1736 dans un article qu'il décrivait comme relevant de
la \og{} géométrie de position \fg{} --- une géométrie sans aucune longueur ni aucun angle, rien que des
connexions. C'est la date de naissance conventionnelle de la théorie des graphes autant que de la
topologie. Euler a démontré la nécessité de sa condition ; la moitié \og{} suffisance \fg{} de la question (b)
n'a été rédigée soigneusement que par Carl Hierholzer en 1873. Le problème récompense ce que la
topologie récompense toujours : trouver la seule quantité qui ne peut pas changer.

\begin{refs}
\item Contexte : Problème des sept ponts de Königsberg --- Wikipédia (\url{https://fr.wikipedia.org/wiki/Probl%C3%A8me_des_sept_ponts_de_K%C3%B6nigsberg})
\item Contexte : Graphe eulérien (chemin eulérien) --- Wikipédia (\url{https://fr.wikipedia.org/wiki/Graphe_eul%C3%A9rien})
\item Lecture : M. A. Armstrong, \textit{Basic Topology}, ch. 1
\end{refs}

\bigskip

\begin{problem}[{La formule des polyèdres d'Euler et les cinq solides de Platon --- Lycée}]
\textbf{TOP-02.} Soit un polyèdre convexe à $ V $ sommets, $ E $ arêtes et $ F $
faces.
\begin{enumerate}
\item[(a)] Démontrez la formule d'Euler
\[ V - E + F = 2 . \]
(Une voie honnête : aplatissez le polyèdre en un réseau plan connexe et raisonnez par récurrence sur
le nombre d'arêtes.)
\item[(b)] Un solide de Platon est un polyèdre convexe dont les faces sont des $ p $-gones réguliers
congruents, avec le même nombre $ q $ de faces se rencontrant en chaque sommet. À l'aide de (a),
démontrez que $ (p, q) $ doit vérifier $ \frac{1}{p} + \frac{1}{q} > \frac{1}{2} $, et
déduisez-en qu'exactement cinq couples sont possibles.
\end{enumerate}

\end{problem}

\noindent Euler a annoncé la formule dans une lettre à Goldbach en 1750, en avouant qu'il ne
savait pas encore la démontrer ; Legendre en a donné une démonstration en 1794 et Cauchy une autre en
1813, et la longue querelle sur ce qui compte exactement comme un \og{} polyèdre \fg{} est devenue le sujet du
dialogue classique de Lakatos, \textit{Proofs and Refutations}. Les cinq solides eux-mêmes closent les
\textit{Éléments} d'Euclide (livre XIII). La quantité $ V - E + F $ est la caractéristique d'Euler --- le
tout premier invariant topologique découvert, et l'ancêtre de tout ce que l'on trouve en TOP-19 et
TOP-20.

\begin{refs}
\item Contexte : Caractéristique d'Euler --- Wikipédia (\url{https://fr.wikipedia.org/wiki/Caract%C3%A9ristique_d'Euler})
\item Contexte : Graphe planaire --- Wikipédia (\url{https://fr.wikipedia.org/wiki/Graphe_planaire})
\item Lecture : M. A. Armstrong, \textit{Basic Topology}, ch. 1
\end{refs}

\bigskip

\begin{problem}[{Expériences avec un ruban de Möbius, rendues rigoureuses --- Lycée}]
\textbf{TOP-03.} Fabriquez un ruban de Möbius : donnez un demi-tour à une bande de papier et collez
les extrémités. Modélisez-le mathématiquement par le carré $ [0,1] \times [0,1] $ dont les bords
gauche et droit sont recollés par $ (0, y) \sim (1, 1-y) $.
\begin{enumerate}
\item[(a)] Démontrez que le bord du ruban de Möbius est une seule courbe fermée (le cylindre ordinaire en a
deux).
\item[(b)] Prédisez ce que l'on obtient en coupant le long de la ligne médiane $ y = \tfrac12 $ ;
démontrez-le ensuite : montrez que la bande coupée est connexe, et déterminez le nombre de demi-tours
qu'elle porte.
\item[(c)] Faites de même pour la coupe le long de la ligne $ y = \tfrac13 $ : démontrez que le résultat a
exactement deux morceaux et décrivez, avec démonstration à partir du modèle de recollement, comment
ils sont joints.
\item[(d)] Montrez qu'une fourmi qui fait un tour complet le long de la ligne médiane revient inversée comme
dans un miroir --- précisez, à l'aide de la règle de recollement, ce que signifie \og{} le ruban n'a qu'une
seule face \fg{}.
\end{enumerate}

\end{problem}

\noindent August Möbius et Johann Listing ont découvert le ruban indépendamment en 1858
(Listing, qui a aussi forgé le mot \og{} topologie \fg{}, a publié le premier). C'est la plus simple des
surfaces non orientables, et les expériences aux ciseaux --- familières des spectacles de magie ---
deviennent des théorèmes dès que le modèle de papier est remplacé par le carré recollé, technique même
qui construit le tore et la bouteille de Klein en TOP-09. La non-orientabilité, rencontrée ici pour la
première fois, est l'un des deux invariants qui classifient toutes les surfaces (TOP-19).

\begin{refs}
\item Contexte : Ruban de Möbius --- Wikipédia (\url{https://fr.wikipedia.org/wiki/Ruban_de_M%C3%B6bius})
\item Lecture : M. A. Armstrong, \textit{Basic Topology}, ch. 1
\end{refs}

\bigskip

\begin{problem}[{Antipodes sur l'équateur --- Lycée}]
\textbf{TOP-04.} Supposez que la température varie continûment le long de
l'équateur terrestre. Démontrez qu'à tout instant donné il existe deux points diamétralement opposés
(antipodaux) sur l'équateur ayant exactement la même température. Formellement : si
$ f : [0, 2\pi] \to \mathbb{R} $ est continue avec $ f(0) = f(2\pi) $, montrez qu'il existe
$ t $ tel que $ f(t) = f(t + \pi) $ (les angles étant pris modulo $ 2\pi $). Vous pouvez
utiliser le théorème des valeurs intermédiaires : une fonction continue qui change de signe sur un
intervalle y admet un zéro.
\end{problem}

\noindent C'est le cas unidimensionnel du théorème de Borsuk--Ulam (TOP-16), et c'est
l'illustration la plus nette possible de la façon dont la topologie démontre l'existence : aucune
formule ne localise jamais les deux points, et pourtant un argument de deux lignes montre qu'ils sont
forcément là. Le théorème des valeurs intermédiaires lui-même a été démontré rigoureusement pour la
première fois par Bolzano en 1817, dans un article explicitement consacré au remplacement de l'évidence
géométrique par la démonstration --- le geste fondateur de l'analyse comme de la topologie
ensembliste.

\begin{refs}
\item Contexte : Théorème des valeurs intermédiaires --- Wikipédia (\url{https://fr.wikipedia.org/wiki/Th%C3%A9or%C3%A8me_des_valeurs_interm%C3%A9diaires})
\item Contexte : Théorème de Borsuk-Ulam --- Wikipédia (\url{https://fr.wikipedia.org/wiki/Th%C3%A9or%C3%A8me_de_Borsuk-Ulam})
\end{refs}

\bigskip

\begin{problem}[{Un point fixe sur l'intervalle --- Lycée, short}]
\textbf{TOP-25.} Soit $ f : [0,1] \to [0,1] $ continue. Démontrez qu'il existe
$ x \in [0,1] $ tel que $ f(x) = x $.
\end{problem}

\noindent C'est le théorème du point fixe de Brouwer en dimension un, et toute la
démonstration tient en une ligne dès que l'on regarde $ g(x) = f(x) - x $ aux deux extrémités.
Brouwer a démontré le cas général vers 1910 ; l'énoncé en dimension supérieure (TOP-14) demande de la
machinerie, mais le cas de l'intervalle ne demande que le théorème des valeurs intermédiaires de
Bolzano --- un bon endroit pour sentir comment une hypothèse purement qualitative force un objet à
exister.

\begin{refs}
\item Contexte : Théorème du point fixe de Brouwer --- Wikipédia (\url{https://fr.wikipedia.org/wiki/Th%C3%A9or%C3%A8me_du_point_fixe_de_Brouwer})
\item Contexte : Théorème des valeurs intermédiaires --- Wikipédia (\url{https://fr.wikipedia.org/wiki/Th%C3%A9or%C3%A8me_des_valeurs_interm%C3%A9diaires})
\end{refs}

\bigskip

\begin{problem}[{Dedans et dehors, comptés par parité --- Lycée, short}]
\textbf{TOP-26.} Soit $ P $ un polygone simple fermé du plan (un nombre fini de
segments, se rencontrant seulement en des extrémités communes, sans auto-intersection) et soit $ q $
un point hors de $ P $. Démontrez que le nombre de fois qu'une demi-droite issue de $ q $ traverse
$ P $ est pair pour toute telle demi-droite, ou impair pour toute telle demi-droite --- la parité ne
dépend pas de la demi-droite que vous lancez, pourvu qu'elle ne passe par aucun sommet de
$ P $.
\end{problem}

\noindent Cette parité est exactement ce que signifie \og{} dedans \fg{} pour un polygone, et la
démontrer est le noyau polygonal honnête du théorème de Jordan : Camille Jordan a énoncé le théorème
général en 1887, et Oswald Veblen en a donné en 1905 une démonstration largement tenue pour la première
complète. Le même argument, sous le nom de lancer de rayon, est ce qu'un moteur graphique exécute chaque
fois qu'il décide si votre curseur est à l'intérieur d'une forme.

\begin{refs}
\item Contexte : Théorème de Jordan --- Wikipédia (\url{https://fr.wikipedia.org/wiki/Th%C3%A9or%C3%A8me_de_Jordan})
\item Contexte : Point dans un polygone --- Wikipédia (en anglais) (\url{https://en.wikipedia.org/wiki/Point_in_polygon})
\end{refs}

\bigskip

\begin{problem}[{Eau, gaz, électricité --- Lycée}]
\textbf{TOP-27.} Trois maisons doivent chacune être reliées à trois services --- eau, gaz,
électricité --- par des conduites tracées dans le plan, sans qu'aucune ne se croise. On note
$ K_{3,3} $ ce graphe et $ K_5 $ le graphe à cinq sommets muni de ses dix arêtes. Vous pouvez
utiliser la formule d'Euler $ V - E + F = 2 $ pour un graphe connexe tracé dans le plan sans
croisement, $ F $ comptant la région extérieure.
\begin{enumerate}
\item[(a)] Démontrez qu'un graphe simple connexe planaire avec $ V \ge 3 $ vérifie
\[ E \le 3V - 6 , \]
en comptant les incidences arête--face, et déduisez-en que $ K_5 $ ne peut pas être tracé dans le
plan sans croisement.
\item[(b)] Démontrez que si un tel graphe ne contient de plus aucun triangle, alors $ E \le 2V - 4 $, et
déduisez-en que $ K_{3,3} $ ne peut pas non plus être tracé --- l'énigme n'a pas de solution.
\item[(c)] Tracez $ K_5 $ et $ K_{3,3} $ sur la surface d'un beignet (un tore) sans croisement, et dites
exactement quelle étape de votre argument de comptage y échoue.
\end{enumerate}

\end{problem}

\noindent Henry Dudeney a publié l'énigme \og{} eau, gaz et électricité \fg{} en 1917 et la disait déjà
ancienne ; des générations de solveurs ont cherché le tracé astucieux qui n'existe pas. La démonstration
qu'il n'existe pas est un petit chef-d'\oe{}uvre de la méthode topologique : rien sur les conduites, un
simple décompte de sommets, d'arêtes et de régions. Kuratowski a démontré en 1930 que $ K_5 $ et
$ K_{3,3} $ sont les \textit{seules} obstructions --- un graphe est planaire exactement quand il ne
contient de subdivision ni de l'un ni de l'autre --- et la question (c) est le premier indice que
\og{} planaire \fg{} n'a jamais concerné le graphe seul, mais la surface sur laquelle on le trace, dont la
caractéristique d'Euler (TOP-02) fait le vrai travail.

\begin{refs}
\item Contexte : Énigme des trois maisons --- Wikipédia (\url{https://fr.wikipedia.org/wiki/%C3%89nigme_des_trois_maisons})
\item Contexte : Graphe planaire --- Wikipédia (\url{https://fr.wikipedia.org/wiki/Graphe_planaire}), Théorème de Kuratowski --- Wikipédia (en anglais) (\url{https://en.wikipedia.org/wiki/Kuratowski%27s_theorem})
\item Lecture : R. Diestel, \textit{Graph Theory}, ch. 4
\end{refs}

\bigskip

\begin{problem}[{Deux alpinistes, une seule altitude --- Lycée}]
\textbf{TOP-37.} Deux profils de montagne sont donnés par des fonctions continues
$ f, g : [0,1] \to [0,1] $ avec $ f(0) = g(0) = 0 $, $ f(1) = g(1) = 1 $ et
$ 0 < f(x), g(x) < 1 $ pour $ 0 < x < 1 $. Un alpiniste parcourt le profil $ f $ et
l'autre le profil $ g $ ; une \textit{traversée conjointe} est un couple d'horaires continus
$ x_1, x_2 : [0,1] \to [0,1] $ avec $ x_1(0) = x_2(0) = 0 $, $ x_1(1) = x_2(1) = 1 $, et
\[ f\big(x_1(u)\big) = g\big(x_2(u)\big) \quad \text{pour tout } u \in [0,1] , \]
de sorte que les deux sont à la même altitude à chaque instant.
\begin{enumerate}
\item[(a)] Esquissez deux profils pour lesquels aucune traversée conjointe ne permet aux deux alpinistes de
n'aller que vers l'est --- à un moment donné, l'un d'eux doit redescendre.
\item[(b)] Supposez maintenant $ f $ et $ g $ affines par morceaux, avec un nombre fini de sommets et de
creux, et telles qu'aucune valeur de maximum ou de minimum local de $ f $ n'égale une valeur de
maximum ou de minimum local de $ g $. Posons
\[ S = \{ (x,y) \in [0,1] \times [0,1] : f(x) = g(y) \} . \]
Démontrez que $ S $ est une réunion finie de segments de droite qui s'assemblent en un nombre fini
d'arcs et de boucles fermées, et que les seules extrémités d'arcs sont les deux coins $ (0,0) $ et
$ (1,1) $.
\item[(c)] Déduisez-en que le morceau de $ S $ contenant $ (0,0) $ est un arc aboutissant à
$ (1,1) $, et lisez une traversée conjointe sur un paramétrage de cet arc --- les alpinistes peuvent
donc toujours le faire.
\item[(d)] Montrez que la continuité seule ne suffit pas : décrivez des profils pour lesquels $ f $ est
constante sur un intervalle tandis que $ g $ oscille une infinité de fois à cette même hauteur, et
expliquez pourquoi le second alpiniste devrait osciller indéfiniment.
\end{enumerate}

\end{problem}

\noindent Tatsuo Homma a résolu une version de ce problème en 1952 et James Whittaker l'a nommé
et posé sous sa forme actuelle en 1966 ; John Philip Huneke a démontré en 1969 qu'une traversée existe
dès qu'aucun des deux profils n'est constant sur un intervalle, et le contre-exemple esquissé en (d)
montre que cette hypothèse ne peut pas être simplement abandonnée. L'exposé de Goodman, Pach et Yap de
1989 a reçu le prix Lester R. Ford de la Mathematical Association of America et a relié l'énigme à la
planification de mouvement des robots et au problème du carré inscrit. La méthode est celle que la
topologie emploie partout : cessez de raisonner sur les alpinistes et raisonnez sur la forme de
l'ensemble de toutes les positions licites.

\begin{refs}
\item Contexte : Problème de l'escalade en montagne --- Wikipédia (en anglais) (\url{https://en.wikipedia.org/wiki/Mountain_climbing_problem})
\item Contexte : Connexité (mathématiques) --- Wikipédia (\url{https://fr.wikipedia.org/wiki/Connexit%C3%A9_(math%C3%A9matiques)})
\item Article : J. E. Goodman, J. Pach, C.-K. Yap, \og{} Mountain climbing, ladder moving, and the ring-width of a polygon \fg{} (1989), \textit{American Mathematical Monthly} 96, 494--510
\end{refs}

\bigskip

\begin{problem}[{Le polyèdre qui casse la formule d'Euler --- Lycée, short}]
\textbf{TOP-38.} Construisez un cadre de tableau polyédrique --- un bloc
rectangulaire percé de part en part d'un trou rectangulaire --- découpé de sorte que chaque face soit un
polygone plat sans trou, puis comptez $ V $, $ E $ et $ F $ et vérifiez que
\[ V - E + F = 0 , \quad \text{et non } 2 . \]
Identifiez exactement quelle étape de la démonstration usuelle de la formule d'Euler (TOP-02) échoue
pour ce solide, et dites ce que le nombre $ V - E + F $ mesure à la place.
\end{problem}

\noindent Simon L'Huilier a publié précisément de telles exceptions en 1812--13, montrant qu'un
solide percé de $ g $ tunnels vérifie $ V - E + F = 2 - 2g $ ; Hessel a ajouté d'autres \og{} monstres \fg{}
en 1832. Imre Lakatos a transformé cette longue querelle --- sur la question de savoir si de tels objets
sont des polyèdres, ou si la définition doit être réparée pour les exclure --- en
\textit{Proofs and Refutations}, le meilleur livre jamais écrit sur la manière dont les mathématiques se
développent réellement. Notez que si vous autorisez le dessus et le dessous à être des faces annulaires
d'un seul tenant, vous retrouvez $ 2 $, et c'est exactement le genre d'échappatoire dont on discutait
au dix-neuvième siècle.

\begin{refs}
\item Contexte : Polyèdre toroïdal --- Wikipédia (en anglais) (\url{https://en.wikipedia.org/wiki/Toroidal_polyhedron})
\item Contexte : Caractéristique d'Euler --- Wikipédia (\url{https://fr.wikipedia.org/wiki/Caract%C3%A9ristique_d'Euler})
\item Lecture : I. Lakatos, \textit{Proofs and Refutations}
\end{refs}

\bigskip


\section*{Licence}

\begin{problem}[{Des distances aux topologies --- Licence}]
\textbf{TOP-05.} Une distance sur un ensemble $ X $ est une fonction
$ d : X \times X \to [0,\infty) $ nulle exactement sur la diagonale, symétrique et vérifiant
l'inégalité triangulaire ; une topologie sur $ X $ est une collection de parties dites \og{} ouvertes \fg{}
contenant $ \varnothing $ et $ X $ et stable par réunion quelconque et par intersection
finie.
\begin{enumerate}
\item[(a)] Démontrez que les parties engendrées par les boules ouvertes d'un espace métrique forment une
topologie.
\item[(b)] Montrez que la distance euclidienne, la distance du taxi
$ d_1(x,y) = |x_1 - y_1| + |x_2 - y_2| $ et la distance du max sur $ \mathbb{R}^2 $ engendrent
toutes la même topologie, tandis que la distance discrète (distance $ 1 $ entre points distincts) en
engendre une autre.
\item[(c)] Montrez que $ d'(x,y) = \frac{|x-y|}{1+|x-y|} $ est une distance sur $\mathbb{R}$ engendrant la topologie
usuelle, bien que $ d' $ soit bornée --- la distance porte donc strictement plus d'information que la
topologie.
\item[(d)] Donnez un exemple d'espace topologique dont la topologie ne provient d'aucune distance, avec
démonstration. (Réfléchissez au nombre de points qu'un espace métrique peut échouer à séparer.)
\end{enumerate}

\end{problem}

\noindent Les axiomes d'espace topologique ont été distillés par Hausdorff dans ses
\textit{Grundzüge der Mengenlehre} (1914) à partir de la théorie des espaces métriques de Fréchet (1906),
en gardant exactement ce qu'il faut pour parler de limites et de continuité et en jetant la distance. Ce
problème est le franchissement rituel de ce seuil : voir quelles distinctions (la bornitude, les
distances précises) s'évaporent, et quels phénomènes nouveaux (les espaces non métrisables) apparaissent
de l'autre côté.

\begin{refs}
\item Contexte : Espace métrique --- Wikipédia (\url{https://fr.wikipedia.org/wiki/Espace_m%C3%A9trique})
\item Contexte : Espace topologique --- Wikipédia (\url{https://fr.wikipedia.org/wiki/Espace_topologique})
\item Lecture : J. R. Munkres, \textit{Topology}, ch. 2
\item Cours : MIT OCW 18.901 Introduction to Topology (en anglais) (\url{https://ocw.mit.edu/courses/18-901-introduction-to-topology-fall-2004/})
\end{refs}

\bigskip

\begin{problem}[{Ouvert, fermé, les deux, ni l'un ni l'autre --- et les quatorze ensembles de Kuratowski --- Licence}]
\textbf{TOP-06.} \og{} Ouvert \fg{} et \og{} fermé \fg{} ne sont pas des contraires ; ce problème calibre les deux
mots.
\begin{enumerate}
\item[(a)] Dans $\mathbb{R}$ muni de la topologie usuelle, donnez des exemples de parties ouvertes et non fermées,
fermées et non ouvertes, à la fois ouvertes et fermées, et ni l'une ni l'autre ; démontrez que $\mathbb{Q}$ n'est
ni l'un ni l'autre.
\item[(b)] L'ouverture est relative : montrez que $ [0, \tfrac12) $ est ouvert dans le sous-espace
$ [0,1] $ mais pas dans $\mathbb{R}$, et formulez précisément ce que signifie \og{} ouvert dans un
sous-espace \fg{}.
\item[(c)] Le problème adhérence--complémentaire de Kuratowski : démontrez qu'en partant d'une seule partie
d'un espace topologique et en appliquant l'adhérence et le passage au complémentaire dans n'importe
quel ordre, on ne peut jamais produire plus de $ 14 $ ensembles distincts --- et exhibez une partie de
$\mathbb{R}$ qui atteint les $ 14 $.
\end{enumerate}

\end{problem}

\noindent La question (c) a été publiée par Kazimierz Kuratowski en 1922, dans l'article même
qui axiomatisait la topologie via l'opérateur d'adhérence, et c'est un rite de passage adoré depuis : la
borne vient d'un petit monoïde d'opérateurs, l'exemple vient de savoir exactement à quel point une partie
de $\mathbb{R}$ peut être pathologique. Les questions (a) et (b) s'attaquent aux deux idées fausses les plus
répandues chez les débutants du domaine --- que \og{} non ouvert \fg{} devrait signifier \og{} fermé \fg{}, et que
l'ouverture est une propriété d'un ensemble plutôt que d'un ensemble dans un espace.

\begin{refs}
\item Contexte : Ouvert (topologie) --- Wikipédia (\url{https://fr.wikipedia.org/wiki/Ouvert_(topologie)}), Fermé (topologie) --- Wikipédia (\url{https://fr.wikipedia.org/wiki/Ferm%C3%A9_(topologie)})
\item Lecture : J. R. Munkres, \textit{Topology}, ch. 2
\end{refs}

\bigskip

\begin{problem}[{Compacité : Borel--Lebesgue et ses limites --- Licence}]
\textbf{TOP-07.} Un espace est \textit{compact} si tout recouvrement par des ouverts admet un
sous-recouvrement fini, et \textit{séquentiellement compact} si toute suite admet une sous-suite
convergente.
\begin{enumerate}
\item[(a)] Démontrez directement à partir de la définition par recouvrements que $ [0,1] $ est
compact.
\item[(b)] Démontrez le théorème de Borel--Lebesgue (Heine--Borel) : une partie de $ \mathbb{R}^n $ est
compacte si et seulement si elle est fermée et bornée.
\item[(c)] Démontrez que dans un espace métrique, compacité et compacité séquentielle sont
équivalentes.
\item[(d)] Montrez que la boule unité fermée de l'espace de suites $ \ell^2 $ est fermée et bornée mais non
compacte --- Borel--Lebesgue échoue en dimension infinie.
\end{enumerate}

\end{problem}

\noindent La compacité est le substitut de la finitude chez le topologue, et il a fallu un
demi-siècle pour trouver la bonne définition : Borel a démontré l'énoncé pour les recouvrements
dénombrables en 1895, Lebesgue et d'autres l'ont généralisé, et la formulation par recouvrements ouverts
a été consacrée par Alexandroff et Urysohn dans les années 1920. La question (d) est le fait qui fait de
l'analyse fonctionnelle une discipline différente de l'analyse en dimension finie --- dans les espaces de
fonctions, les parties fermées bornées ne portent aucune compacité, et il faut inventer des substituts
(équicontinuité, topologies faibles). L'équivalence de la question (c) échoue réellement hors des espaces
métriques : voir TOP-13.

\begin{refs}
\item Contexte : Compacité (mathématiques) --- Wikipédia (\url{https://fr.wikipedia.org/wiki/Compacit%C3%A9_(math%C3%A9matiques)})
\item Contexte : Théorème de Borel-Lebesgue (Heine--Borel) --- Wikipédia (\url{https://fr.wikipedia.org/wiki/Th%C3%A9or%C3%A8me_de_Borel-Lebesgue})
\item Lecture : J. R. Munkres, \textit{Topology}, ch. 3
\item Cours : MIT OCW 18.901 Introduction to Topology (en anglais) (\url{https://ocw.mit.edu/courses/18-901-introduction-to-topology-fall-2004/})
\end{refs}

\bigskip

\begin{problem}[{Connexité et la courbe sinus du topologue --- Licence}]
\textbf{TOP-08.} Un espace est \textit{connexe} s'il n'est pas la réunion de deux ouverts disjoints
non vides, et \textit{connexe par arcs} si deux points quelconques sont joints par un chemin
continu.
\begin{enumerate}
\item[(a)] Démontrez que les images continues d'espaces connexes sont connexes, et déduisez-en le théorème
des valeurs intermédiaires.
\item[(b)] Démontrez que les espaces connexes par arcs sont connexes.
\item[(c)] Soit $ S = \{ (x, \sin(1/x)) : 0 < x \le 1 \} \cup (\{0\} \times [-1,1]) $, la courbe sinus
du topologue. Démontrez que $ S $ est connexe mais non connexe par arcs.
\item[(d)] Utilisez la connexité pour démontrer que $\mathbb{R}$ et $ \mathbb{R}^2 $ ne sont pas
homéomorphes.
\end{enumerate}

\end{problem}

\noindent La connexité semble aller de soi jusqu'à ce que des exemples comme (c) --- un ensemble
que l'on peut dessiner mais dans lequel on ne peut pas marcher --- écartent les deux définitions rivales.
La notion rigoureuse a émergé des travaux de Jordan et de Schoenflies et a reçu sa forme moderne de Riesz
et de Hausdorff ; la courbe sinus est devenue la mise en garde standard au début du vingtième siècle. La
question (d) est le premier cas du problème de l'invariance de la dimension, dont la solution générale
par Brouwer en 1911 (via les outils de TOP-20) a créé le premier triomphe de la topologie
algébrique.

\begin{refs}
\item Contexte : Connexité (mathématiques) --- Wikipédia (\url{https://fr.wikipedia.org/wiki/Connexit%C3%A9_(math%C3%A9matiques)})
\item Contexte : Courbe sinus du topologue --- Wikipédia (\url{https://fr.wikipedia.org/wiki/Courbe_sinus_du_topologue})
\item Lecture : J. R. Munkres, \textit{Topology}, ch. 3
\end{refs}

\bigskip

\begin{problem}[{Recollements : tore, bouteille de Klein, plan projectif --- Licence}]
\textbf{TOP-09.} Le quotient d'un espace par une identification porte la topologie la plus fine
rendant la projection continue. Prenez le carré $ [0,1] \times [0,1] $ et recollez : (i)
$ (0,y) \sim (1,y) $ et $ (x,0) \sim (x,1) $ --- le tore $ T^2 $ ; (ii)
$ (0,y) \sim (1, 1-y) $ et $ (x,0) \sim (x,1) $ --- la bouteille de Klein $ K $ ; (iii) toutes les
identifications du bord $ (x, 0) \sim (1-x, 1) $ et $ (0, y) \sim (1, 1-y) $ --- le plan projectif
réel $ \mathbb{RP}^2 $, autrement dit la sphère dont on identifie les points antipodaux.
\begin{enumerate}
\item[(a)] Démontrez que le quotient tore est homéomorphe à $ S^1 \times S^1 $.
\item[(b)] Démontrez que $ \mathbb{RP}^2 $ tel que défini à partir du carré est homéomorphe au quotient
antipodal de $ S^2 $, et que les deux sont compacts.
\item[(c)] Montrez que la bouteille de Klein est la réunion de deux rubans de Möbius recollés le long de
leurs cercles de bord.
\item[(d)] Calculez $ V - E + F $ pour chaque surface à partir de sa présentation carrée, et vérifiez les
valeurs $ 0, 0, 1 $.
\end{enumerate}

\end{problem}

\noindent Les constructions par quotient sont la façon dont la topologie bâtit des espaces
invisibles dans $ \mathbb{R}^3 $ : Klein a décrit sa bouteille en 1882 (il lui faut quatre dimensions
pour éviter l'auto-intersection), et le plan projectif descend de la géométrie projective des peintres de
la Renaissance en passant par Möbius. La notation \og{} carré fléché \fg{} est la comptabilité standard de
TOP-19, où toute surface compacte sera présentée ainsi. Rendre (a) et (b) pleinement rigoureux --- les
bijections continues d'un compact vers un séparé sont des homéomorphismes --- est l'objet de
l'exercice.

\begin{refs}
\item Contexte : Espace quotient (topologie) --- Wikipédia (\url{https://fr.wikipedia.org/wiki/Espace_quotient_(topologie)})
\item Contexte : Bouteille de Klein --- Wikipédia (\url{https://fr.wikipedia.org/wiki/Bouteille_de_Klein}), Plan projectif réel --- Wikipédia (\url{https://fr.wikipedia.org/wiki/Plan_projectif_r%C3%A9el})
\item Lecture : M. A. Armstrong, \textit{Basic Topology}, ch. 4
\item Lecture : J. R. Munkres, \textit{Topology}, ch. 2
\end{refs}

\bigskip

\begin{problem}[{Distinguer les espaces --- Licence}]
\textbf{TOP-10.} Un \textit{invariant topologique} est une propriété préservée par tout
homéomorphisme.
\begin{enumerate}
\item[(a)] Démontrez que la compacité et la connexité sont des invariants topologiques.
\item[(b)] Un \textit{point de coupure} d'un espace connexe est un point dont le retrait le déconnecte.
Démontrez que les homéomorphismes envoient les points de coupure sur des points de coupure.
\item[(c)] À l'aide de (a) et (b), démontrez que parmi les quatre espaces suivants, aucun couple n'est formé
d'espaces homéomorphes : l'intervalle fermé $ [0,1] $, le cercle $ S^1 $, la droite réelle $\mathbb{R}$ et le
huit (deux cercles joints en un point).
\item[(d)] Expliquez pourquoi les arguments de type \og{} nombre de points de coupure \fg{} ne peuvent jamais
distinguer $ \mathbb{R}^2 $ de $ \mathbb{R}^3 $, et nommez l'invariant qui, lui, le fait.
\end{enumerate}

\end{problem}

\noindent C'est le jeu fondamental de la topologie --- pour démontrer que deux espaces sont
\textit{différents} il faut produire un invariant, puisque aucune quantité d'échecs à trouver un
homéomorphisme ne démontre qu'il n'en existe pas. Les arguments par points de coupure sont les invariants
les plus anciens et les plus élémentaires ; leur insuffisance pour la question (d), comprise vers 1900,
est précisément ce qui a poussé Poincaré à attacher des groupes aux espaces (TOP-12) et Brouwer à
développer la théorie du degré. Les quatre espaces de (c) reviendront tout au long de cette page comme cas
de test de base.

\begin{refs}
\item Contexte : Propriété topologique --- Wikipédia (\url{https://fr.wikipedia.org/wiki/Propri%C3%A9t%C3%A9_topologique})
\item Contexte : Homéomorphisme --- Wikipédia (\url{https://fr.wikipedia.org/wiki/Hom%C3%A9omorphisme})
\item Lecture : J. R. Munkres, \textit{Topology}, ch. 3
\end{refs}

\bigskip

\begin{problem}[{Le lemme de Sperner et un point fixe --- Licence}]
\textbf{TOP-11.} Subdivisez un triangle $ ABC $ en petits triangles, et coloriez chaque sommet de
la subdivision avec l'une de trois couleurs, de sorte que $ A, B, C $ reçoivent trois couleurs
différentes et que tout sommet situé sur un côté de $ ABC $ n'utilise que les deux couleurs des
extrémités de ce côté.
\begin{enumerate}
\item[(a)] Démontrez le lemme de Sperner : un petit triangle porte les trois couleurs à ses sommets (et même,
il y en a un nombre impair).
\item[(b)] Utilisez (a) pour démontrer le théorème du point fixe de Brouwer en dimension deux : toute
application continue d'un triangle fermé (autrement dit du disque fermé $ D^2 $) dans lui-même admet
un point fixe.
\item[(c)] Démontrez le cas unidimensionnel --- toute application continue $ f : [0,1] \to [0,1] $ admet un
point fixe --- directement à partir du théorème des valeurs intermédiaires.
\end{enumerate}

\end{problem}

\noindent Emanuel Sperner a démontré son lemme combinatoire en 1928, et en moins d'un an
Knaster, Kuratowski et Mazurkiewicz en avaient déduit le théorème de Brouwer (démontré à l'origine en
1911 avec de lourdes machineries) --- l'une des toutes premières démonstrations qu'un comptage discret peut
porter des conséquences continues. Le lemme est aujourd'hui le moteur des algorithmes de partage équitable
en économie et du calcul des équilibres économiques de Scarf. TOP-14 donne la démonstration par lacets et
les corollaires classiques du théorème.

\begin{refs}
\item Contexte : Lemme de Sperner --- Wikipédia (\url{https://fr.wikipedia.org/wiki/Lemme_de_Sperner})
\item Contexte : Théorème du point fixe de Brouwer --- Wikipédia (\url{https://fr.wikipedia.org/wiki/Th%C3%A9or%C3%A8me_du_point_fixe_de_Brouwer})
\item Lecture : M. A. Armstrong, \textit{Basic Topology}, ch. 5
\end{refs}

\bigskip

\begin{problem}[{Le groupe fondamental, et le lacet qui compte --- Licence}]
\textbf{TOP-12.} Pour un espace $ X $ muni d'un point base $ x_0 $, soit
$ \pi_1(X, x_0) $ l'ensemble des classes d'homotopie de lacets en $ x_0 $, muni de la composition
donnée par la concaténation.
\begin{enumerate}
\item[(a)] Démontrez que c'est un groupe.
\item[(b)] Démontrez qu'une application continue $ f : X \to Y $ induit un morphisme
$ f_* : \pi_1(X, x_0) \to \pi_1(Y, f(x_0)) $, et que des espaces homéomorphes connexes par arcs ont
des groupes fondamentaux isomorphes.
\item[(c)] Démontrez que $ \pi_1 $ d'une partie convexe de $ \mathbb{R}^n $ est trivial.
\item[(d)] Démontrez le premier vrai calcul : $ \pi_1(S^1) \cong \mathbb{Z} $. Utilisez l'application
$ p(t) = (\cos 2\pi t, \sin 2\pi t) $ de $\mathbb{R}$ vers le cercle : montrez que tout lacet se relève de
manière unique en un chemin de $\mathbb{R}$, que les homotopies se relèvent également, et que l'extrémité du
relèvement --- l'indice --- est un isomorphisme bien défini sur $\mathbb{Z}$.
\end{enumerate}

\end{problem}

\noindent Poincaré a introduit le groupe fondamental dans \textit{Analysis Situs} (1895),
l'article qui a fondé la topologie algébrique ; c'était la première fois qu'un groupe était attaché à un
espace. Le calcul de la question (d) est le prototype de la théorie des revêtements (TOP-17) et le moteur
derrière le théorème de non-rétraction, le théorème de Brouwer et Borsuk--Ulam en dimension deux.
Comprendre cette seule démonstration à fond --- ce qui se relève, pourquoi les relèvements sont uniques, où
intervient la complétude de $\mathbb{R}$ --- vaut un semestre de lecture passive.

\begin{refs}
\item Contexte : Groupe fondamental --- Wikipédia (\url{https://fr.wikipedia.org/wiki/Groupe_fondamental}), Homotopie --- Wikipédia (\url{https://fr.wikipedia.org/wiki/Homotopie})
\item Contexte : Indice d'un lacet --- Wikipédia (\url{https://fr.wikipedia.org/wiki/Indice_(analyse_complexe)})
\item Lecture : A. Hatcher, \textit{Algebraic Topology}, ch. 1 --- gratuit sur la page de l'auteur (en anglais) (\url{https://pi.math.cornell.edu/~hatcher/AT/ATpage.html})
\item Lecture : J. R. Munkres, \textit{Topology}, ch. 9
\end{refs}

\bigskip

\begin{problem}[{Une application quotient qui n'est ni ouverte ni fermée --- Licence, short}]
\textbf{TOP-28.} Une surjection $ q : X \to Y $ est une \textit{application
quotient} si une partie $ U \subseteq Y $ est ouverte exactement quand $ q^{-1}(U) $ est ouverte
dans $ X $. Produisez une application quotient qui n'est ni ouverte ni fermée, et démontrez les trois
affirmations à propos de votre exemple.
\end{problem}

\noindent Les applications quotients sont la façon dont tout espace recollé de cette page est
construit --- le tore, la bouteille de Klein, le plan projectif de TOP-09 --- aussi importe-t-il que la
définition soit un énoncé sur les parties ouvertes en aval, et non une promesse que les images d'ouverts
sont ouvertes. Les étudiants supposent couramment cette promesse ; l'objet de cet exercice est de rompre
l'habitude une bonne fois, un exemple en main.

\begin{refs}
\item Contexte : Espace quotient (topologie) --- Wikipédia (\url{https://fr.wikipedia.org/wiki/Espace_quotient_(topologie)})
\item Lecture : J. R. Munkres, \textit{Topology}, \S{}22
\end{refs}

\bigskip

\begin{problem}[{Le lemme du nombre de Lebesgue --- Licence, short}]
\textbf{TOP-29.} Soit $ (X,d) $ un espace métrique compact et soit
$ \mathcal{U} $ un recouvrement ouvert de $ X $. Démontrez qu'il existe $ \delta > 0 $ --- un
\textit{nombre de Lebesgue} pour $ \mathcal{U} $ --- tel que toute partie de $ X $ de diamètre
strictement inférieur à $ \delta $ soit contenue dans un seul membre de $ \mathcal{U} $.
\end{problem}

\noindent Nommé d'après Henri Lebesgue, ce petit lemme est la charnière sur laquelle tourne une
quantité surprenante de choses : il donne la démonstration en une ligne du fait qu'une fonction continue
sur un espace métrique compact est uniformément continue, et c'est lui qui permet de découper un chemin
ou un carré d'homotopie en morceaux assez petits pour tenir dans un seul ouvert --- le premier geste des
démonstrations du théorème de relèvement des chemins (TOP-17) et de Seifert--van Kampen (TOP-18).

\begin{refs}
\item Contexte : Nombre de Lebesgue --- Wikipédia (\url{https://fr.wikipedia.org/wiki/Nombre_de_Lebesgue})
\item Lecture : J. R. Munkres, \textit{Topology}, \S{}27
\end{refs}

\bigskip

\begin{problem}[{Topologie produit contre topologie des boîtes --- Licence}]
\textbf{TOP-30.} Sur un produit infini $ \prod_{\alpha \in A} X_\alpha $ il existe deux topologies
naturelles candidates. La topologie \textit{des boîtes} a pour base tous les produits
$ \prod_\alpha U_\alpha $ avec chaque $ U_\alpha $ ouvert ; la topologie \textit{produit} ne retient
que les boîtes pour lesquelles $ U_\alpha = X_\alpha $ sauf pour un nombre fini de $ \alpha $. On
note $ \mathbb{R}^\omega $ le produit d'une infinité dénombrable de copies de $\mathbb{R}$.
\begin{enumerate}
\item[(a)] Montrez que sur $ \mathbb{R}^\omega $ la topologie des boîtes est strictement plus fine que la
topologie produit.
\item[(b)] Démontrez qu'une application $ f : Z \to \prod_\alpha X_\alpha $ est continue pour la topologie
produit si et seulement si chaque fonction coordonnée $ \pi_\alpha \circ f $ est continue, et montrez
que cela échoue pour la topologie des boîtes en examinant $ t \mapsto (t, t, t, \ldots) $ de $\mathbb{R}$ vers
$ \mathbb{R}^\omega $.
\item[(c)] Démontrez que $ \mathbb{R}^\omega $ est connexe pour la topologie produit mais pas pour la
topologie des boîtes.
\item[(d)] Montrez que $ \{0,1\}^\omega $ porte la topologie discrète pour la topologie des boîtes, et
concluez que le théorème de Tykhonov (TOP-13) est faux pour les boîtes. Énoncez et démontrez ensuite la
propriété universelle qui singularise la topologie produit : c'est la topologie la moins fine rendant
continue chaque projection $ \pi_\alpha $.
\end{enumerate}

\end{problem}

\noindent Quand Tykhonov a introduit les produits infinis en 1930 il a dû choisir, et le choix
paraît d'abord pervers --- pourquoi refuser d'appeler ouverte une boîte manifeste ? La question (b) est la
réponse : seule la topologie la moins fine rend les applications \textit{vers} un produit aussi faciles à
construire que les coordonnées le suggèrent, ce qui est précisément la propriété universelle qu'un produit
catégorique doit avoir, et seule la topologie la moins fine conserve la compacité. La topologie des boîtes
survit comme contre-exemple standard, le lieu où meurt discrètement tout théorème commode sur les
produits.

\begin{refs}
\item Contexte : Topologie produit --- Wikipédia (\url{https://fr.wikipedia.org/wiki/Topologie_produit})
\item Contexte : Topologie des boîtes --- Wikipédia (\url{https://fr.wikipedia.org/wiki/Topologie_des_bo%C3%AEtes})
\item Lecture : J. R. Munkres, \textit{Topology}, \S{}19--20
\item Cours : MIT OCW 18.901 Introduction to Topology (en anglais) (\url{https://ocw.mit.edu/courses/18-901-introduction-to-topology-fall-2004/})
\end{refs}

\bigskip

\begin{problem}[{Une courbe qui remplit un carré --- Licence}]
\textbf{TOP-39.} Une \textit{courbe remplissante} est une surjection continue
$ f : [0,1] \to [0,1] \times [0,1] $. Construisez-en une, puis voyez ce qu'elle ne peut pas
être.
\begin{enumerate}
\item[(a)] Construisez la courbe de Hilbert comme limite : définissez récursivement des applications
polygonales $ f_n : [0,1] \to [0,1]^2 $, chacune visitant tour à tour les quatre quadrants du carré,
et démontrez que $ (f_n) $ converge uniformément, de sorte que la limite $ f $ est continue.
\item[(b)] Démontrez que $ f $ est surjective. (L'image d'un compact est compacte, donc fermée ; il suffit
donc de montrer que l'image est dense.)
\item[(c)] Démontrez qu'aucune courbe remplissante ne peut être injective : montrez qu'une bijection continue
d'un compact sur un espace séparé est un homéomorphisme, puis que $ [0,1] $ et $ [0,1]^2 $ ne sont
pas homéomorphes (TOP-10 fournit l'invariant).
\item[(d)] Expliquez pourquoi (a) et (c) réunies ne contredisent pas l'invariance de la dimension de Brouwer,
et dites précisément quelle hypothèse de ce théorème la courbe met en défaut.
\end{enumerate}

\end{problem}

\noindent Cantor a montré en 1878 que $ [0,1] $ et $ [0,1]^2 $ ont le même cardinal, ce qui
était déjà un scandale ; Netto a démontré la même année qu'aucune telle bijection ne peut être continue.
Cela laissait une question ouverte, et Peano y a répondu en 1890 par une construction purement
arithmétique --- son article ne contient aucune figure --- d'une courbe continue passant par tout point d'un
carré ; Hilbert en a donné la version géométrique un an plus tard, et c'est celle que tout le monde
dessine. La morale est que \og{} courbe \fg{}, définie comme \og{} image continue d'un intervalle \fg{}, est une
définition beaucoup trop généreuse, et la réparer a occupé les topologues pendant des décennies. Ces mêmes
courbes indexent aujourd'hui des bases de données et des tuiles d'images, parce que des points proches sur
la droite restent proches dans le carré.

\begin{refs}
\item Contexte : Courbe remplissante --- Wikipédia (\url{https://fr.wikipedia.org/wiki/Courbe_remplissante})
\item Contexte : Courbe de Hilbert --- Wikipédia (\url{https://fr.wikipedia.org/wiki/Courbe_de_Hilbert}), Courbe de Peano --- Wikipédia (\url{https://fr.wikipedia.org/wiki/Courbe_de_Peano})
\item Lecture : J. R. Munkres, \textit{Topology}, ch. 7
\item Lecture : H. Sagan, \textit{Space-Filling Curves} (Springer, 1994)
\end{refs}

\bigskip

\begin{problem}[{Même type d'homotopie, espace différent --- Licence, short}]
\textbf{TOP-40.} Démontrez que le ruban de Möbius se rétracte par déformation sur
son cercle central, de sorte qu'il est homotopiquement équivalent à $ S^1 $ et que son groupe
fondamental est $ \mathbb{Z} $. Démontrez ensuite que le ruban n'est néanmoins \textit{pas} homéomorphe
à $ S^1 $.
\end{problem}

\noindent L'équivalence d'homotopie, isolée dans les années 1930 dans les travaux de Borsuk et
de Hurewicz, est la relation plus grossière que voient réellement tous les invariants algébriques de cette
page : le groupe fondamental de TOP-12 et l'homologie de TOP-20 ne peuvent pas distinguer le ruban de son
cercle central, parce que rien de bâti sur la déformation continue ne le peut. Les distinguer demande un
invariant véritablement topologique, et l'invariant élémentaire de TOP-10 le fait en une ligne. Savoir à
quelles questions votre invariant est aveugle, c'est la moitié de son usage.

\begin{refs}
\item Contexte : Homotopie --- Wikipédia (\url{https://fr.wikipedia.org/wiki/Homotopie}), Rétraction (topologie) --- Wikipédia (\url{https://fr.wikipedia.org/wiki/R%C3%A9traction})
\item Contexte : Ruban de Möbius --- Wikipédia (\url{https://fr.wikipedia.org/wiki/Ruban_de_M%C3%B6bius})
\item Lecture : A. Hatcher, \textit{Algebraic Topology}, ch. 0 --- gratuit sur la page de l'auteur (en anglais) (\url{https://pi.math.cornell.edu/~hatcher/AT/ATpage.html})
\end{refs}

\bigskip

\begin{problem}[{Construire des espaces avec des cellules --- Licence}]
\textbf{TOP-41.} Un \textit{CW-complexe} fini se construit à partir d'un ensemble fini de points en
attachant des cellules : à l'étape $ n $ on recolle des $ n $-boules fermées sur la partie déjà
bâtie, le long d'applications continues de leurs sphères de bord. Si $ c_n $ est le nombre de
$ n $-cellules, la caractéristique d'Euler est
\[ \chi(X) = \sum_{n \ge 0} (-1)^n c_n . \]
\begin{enumerate}
\item[(a)] Donnez des structures CW ayant le moins de cellules possible sur : la sphère $ S^n $ (deux
cellules suffisent) ; le tore $ T^2 $ ; le bouquet de $ k $ cercles ; la surface fermée orientable
de genre $ g $ ; et $ \mathbb{RP}^n $ (une cellule en chaque dimension $ 0, 1, \ldots, n $). Dans
chaque cas, dites exactement quelles sont les applications d'attachement.
\item[(b)] Démontrez qu'une partie compacte d'un CW-complexe ne rencontre qu'un nombre fini de cellules
ouvertes, et déduisez-en qu'un CW-complexe est compact exactement quand il a un nombre fini de
cellules.
\item[(c)] Calculez $ \chi $ pour tous les espaces de (a). Vérifiez $ \chi = 2 - 2g $ pour la surface de
genre $ g $, et que $ \chi(\mathbb{RP}^n) $ vaut $ 1 $ pour $ n $ pair et $ 0 $ pour $ n $
impair.
\item[(d)] Démontrez que pour les CW-complexes finis
\[ \chi(X \vee Y) = \chi(X) + \chi(Y) - 1 , \qquad \chi(X \times Y) = \chi(X)\,\chi(Y) , \]
et utilisez ces formules pour calculer $ \chi $ du tore de dimension $ n $, $ (S^1)^n $.
\end{enumerate}

\end{problem}

\noindent J. H. C. Whitehead a introduit les CW-complexes dans \og{} Combinatorial homotopy I \fg{}
(1949) --- les lettres signifient \textit{closure-finite} (fermeture finie) avec la topologie \textit{weak}
(faible) --- et ils ont supplanté presque aussitôt les complexes simpliciaux comme catégorie de travail de
la topologie algébrique. La raison en est l'économie, visible dès (a) : la sphère $ S^n $ exige au moins
$ n + 2 $ sommets comme complexe simplicial mais seulement deux cellules ici, et $ \mathbb{RP}^n $
devient un objet à une cellule par dimension dont on lit presque directement l'homologie. Tout ce qui suit
sur cette page --- les calculs de Mayer--Vietoris de TOP-43, la théorie de Lefschetz de TOP-44, la propre
question ouverte de Whitehead en TOP-34 --- est énoncé pour des espaces construits ainsi.

\begin{refs}
\item Contexte : CW-complexe --- Wikipédia (\url{https://fr.wikipedia.org/wiki/CW-complexe})
\item Contexte : J. H. C. Whitehead --- Wikipédia (\url{https://fr.wikipedia.org/wiki/J._H._C._Whitehead})
\item Lecture : A. Hatcher, \textit{Algebraic Topology}, ch. 0 et appendice --- gratuit sur la page de l'auteur (en anglais) (\url{https://pi.math.cornell.edu/~hatcher/AT/ATpage.html})
\item Cours : MIT OCW 18.905 Algebraic Topology I (en anglais) (\url{https://ocw.mit.edu/courses/18-905-algebraic-topology-i-fall-2016/})
\end{refs}

\bigskip


\section*{Master}

\begin{problem}[{La compacité au-delà des espaces métriques --- Master}]
\textbf{TOP-13.} TOP-07 a montré que compacité par recouvrements et compacité séquentielle
coïncident dans les espaces métriques. Montrez qu'en général elles sont indépendantes :
\begin{enumerate}
\item[(a)] Démontrez le théorème de Tykhonov pour un produit quelconque de copies de $ [0,1] $ (tout produit
d'espaces compacts est compact pour la topologie produit), et utilisez-le pour montrer que le cube
$ [0,1]^{[0,1]} $ --- toutes les applications de $ [0,1] $ dans lui-même, avec la convergence
ponctuelle --- est compact mais \textit{non} séquentiellement compact.
\item[(b)] Démontrez que l'espace des ordinaux dénombrables $ [0, \omega_1) $ muni de la topologie de
l'ordre est séquentiellement compact mais non compact.
\item[(c)] Localisez l'échec dans chaque direction : quelle étape de la démonstration valable en espace
métrique se casse, et pourquoi ?
\end{enumerate}

\end{problem}

\noindent Tykhonov a démontré son théorème en 1930 pour les cubes et énoncé le cas général en
1935 ; Kakutani et d'autres ont montré plus tard qu'il équivaut à l'axiome du choix, et Kelley l'a appelé
\og{} probablement le plus important théorème isolé de la topologie générale \fg{} --- il sous-tend le théorème de
Banach--Alaoglu et le théorème de compacité de la logique. Les deux contre-exemples de (a) et (b) sont les
témoins standard du fait que les suites ne suffisent pas à détecter la topologie en général, ce qui est la
raison de l'invention des filtres et des suites généralisées (Moore--Smith 1922, Cartan 1937).

\begin{refs}
\item Contexte : Théorème de Tykhonov --- Wikipédia (\url{https://fr.wikipedia.org/wiki/Th%C3%A9or%C3%A8me_de_Tykhonov})
\item Contexte : Compacité séquentielle --- Wikipédia (\url{https://fr.wikipedia.org/wiki/Compacit%C3%A9_s%C3%A9quentielle}), Topologie produit --- Wikipédia (\url{https://fr.wikipedia.org/wiki/Topologie_produit})
\item Lecture : J. R. Munkres, \textit{Topology}, ch. 5
\end{refs}

\bigskip

\begin{problem}[{Pas de rétraction, points fixes, et le théorème fondamental de l'algèbre --- Master}]
\textbf{TOP-14.} On admet le calcul $ \pi_1(S^1) \cong \mathbb{Z} $ de TOP-12.
\begin{enumerate}
\item[(a)] À l'aide de $ \pi_1(S^1) \cong \mathbb{Z} $ (TOP-12), démontrez le théorème de non-rétraction :
il n'existe pas d'application continue $ r : D^2 \to S^1 $ se restreignant à l'identité sur le cercle
de bord.
\item[(b)] Déduisez-en le théorème du point fixe de Brouwer pour $ D^2 $, et démontrez réciproquement que le
théorème du point fixe entraîne la non-rétraction --- les deux énoncés sont équivalents.
\item[(c)] Utilisez les indices de lacets pour démontrer le théorème fondamental de l'algèbre : tout polynôme
complexe non constant a une racine dans $\mathbb{C}$.
\item[(d)] Énoncez les théorèmes de non-rétraction et de point fixe pour $ D^n $, $ n \ge 3 $, et
expliquez précisément pourquoi $ \pi_1 $ ne peut pas les démontrer --- identifiez quelle propriété de
$ S^{n-1} $ fait défaut, et quel invariant (TOP-20) s'impose à la place.
\end{enumerate}

\end{problem}

\noindent Brouwer a démontré son théorème du point fixe en 1911, dans la même salve de travaux
qui a établi l'invariance de la dimension ; l'équivalence avec la non-rétraction est l'exemple le plus net
du geste caractéristique de la topologie, qui convertit un énoncé sur toutes les applications en un énoncé
sur l'existence d'une seule. La démonstration du théorème fondamental de l'algèbre en (c) est pour
l'essentiel la première démonstration de Gauss (1799) rendue rigoureuse --- l'indice fournit exactement
l'argument de continuité sur lequel Gauss passait la main.

\begin{refs}
\item Contexte : Théorème du point fixe de Brouwer --- Wikipédia (\url{https://fr.wikipedia.org/wiki/Th%C3%A9or%C3%A8me_du_point_fixe_de_Brouwer})
\item Contexte : Rétraction (topologie) --- Wikipédia (\url{https://fr.wikipedia.org/wiki/R%C3%A9traction})
\item Lecture : A. Hatcher, \textit{Algebraic Topology}, ch. 1 --- gratuit sur la page de l'auteur (en anglais) (\url{https://pi.math.cornell.edu/~hatcher/AT/ATpage.html})
\end{refs}

\bigskip

\begin{problem}[{Le théorème de la boule chevelue --- Master}]
\textbf{TOP-15.} Un champ de vecteurs tangent sur la sphère $ S^2 $ associe continûment à chaque
point $ x $ un vecteur $ v(x) $ perpendiculaire à $ x $.
\begin{enumerate}
\item[(a)] Démontrez le théorème de la boule chevelue : tout champ de vecteurs tangent continu sur $ S^2 $
s'annule quelque part. (Une voie : montrez qu'un champ ne s'annulant nulle part fournirait une homotopie
entre l'identité de $ S^2 $ et l'application antipodale, et démontrez que celles-ci ne peuvent pas
être homotopes en comparant les degrés.)
\item[(b)] Montrez par une construction explicite que le théorème échoue sur le tore : exhibez un champ
tangent continu ne s'annulant nulle part.
\item[(c)] Déduisez-en qu'à tout instant il existe un point de la Terre où la vitesse horizontale du vent est
nulle, et démontrez honnêtement ce corollaire météorologique à partir de (a).
\end{enumerate}

\end{problem}

\noindent Poincaré a démontré le cas de dimension deux en 1885 dans le cadre de son théorème de
l'indice pour les champs de vecteurs sur les surfaces --- l'indice total égale la caractéristique d'Euler,
ce qui explique que le tore ($ \chi = 0 $) y échappe ; Brouwer a étendu l'énoncé à toutes les sphères de
dimension paire en 1912. La démonstration analytique d'une page de Milnor (1978) est une démonstration
célèbre du fait que même des théorèmes centenaires peuvent être redémontrés à partir de rien. Le résultat
sème toute la théorie des classes caractéristiques : des obstructions aux champs et aux sections mesurées
par la topologie.

\begin{refs}
\item Contexte : Théorème de la boule chevelue --- Wikipédia (\url{https://fr.wikipedia.org/wiki/Th%C3%A9or%C3%A8me_de_la_boule_chevelue})
\item Lecture : A. Hatcher, \textit{Algebraic Topology}, ch. 2 --- gratuit sur la page de l'auteur (en anglais) (\url{https://pi.math.cornell.edu/~hatcher/AT/ATpage.html})
\item Article : J. Milnor, \og{} Analytic proofs of the `hairy ball theorem' and the Brouwer fixed point theorem \fg{} (1978), \textit{American Mathematical Monthly} 85
\end{refs}

\bigskip

\begin{problem}[{Borsuk--Ulam et le sandwich au jambon --- Master}]
\textbf{TOP-16.} Ce que toute application de la sphère vers le plan doit faire subir à un couple de
points antipodaux.
\begin{enumerate}
\item[(a)] Démontrez le théorème de Borsuk--Ulam en dimension deux : toute application continue
$ f : S^2 \to \mathbb{R}^2 $ envoie un couple de points antipodaux sur la même valeur,
$ f(x) = f(-x) $. (Une voie standard utilise la technologie de relèvement des lacets de TOP-12
appliquée à l'équateur.)
\item[(b)] Déduisez-en qu'aucune application continue $ S^2 \to S^1 $ ne commute avec les applications
antipodales, et que $ S^2 $ ne peut pas se plonger dans $ \mathbb{R}^2 $.
\item[(c)] Déduisez-en le théorème du sandwich au jambon : pour trois solides mesurables bornés quelconques de
$ \mathbb{R}^3 $ (pain, jambon, pain), un même plan coupe les trois exactement en deux parts de volume
égal.
\item[(d)] Déduisez-en également la vedette météorologique : à tout instant il existe deux points antipodaux
de la Terre ayant la même température \textit{et} la même pression.
\end{enumerate}

\end{problem}

\noindent Stanis\l{}aw Ulam a posé la conjecture et Karol Borsuk l'a démontrée en 1933 ; le
corollaire du sandwich au jambon avait été posé par Steinhaus et a été démontré via Borsuk--Ulam par Stone
et Tukey en 1942 dans toute sa généralité ($ n $ solides dans $ \mathbb{R}^n $). Le théorème est
strictement plus fort que le théorème du point fixe de Brouwer, et il est devenu ces dernières décennies
un cheval de trait de la combinatoire : la démonstration par Lovász de la conjecture de Kneser (1978) et
le théorème du partage de collier en découlent tous deux, histoire magnifiquement racontée dans le
\textit{Using the Borsuk--Ulam Theorem} de Matou\v{s}ek.

\begin{refs}
\item Contexte : Théorème de Borsuk-Ulam --- Wikipédia (\url{https://fr.wikipedia.org/wiki/Th%C3%A9or%C3%A8me_de_Borsuk-Ulam})
\item Contexte : Théorème du sandwich au jambon --- Wikipédia (\url{https://fr.wikipedia.org/wiki/Th%C3%A9or%C3%A8me_du_sandwich_au_jambon})
\item Lecture : A. Hatcher, \textit{Algebraic Topology}, ch. 1 --- gratuit sur la page de l'auteur (en anglais) (\url{https://pi.math.cornell.edu/~hatcher/AT/ATpage.html})
\end{refs}

\bigskip

\begin{problem}[{Revêtements --- Master}]
\textbf{TOP-17.} Un revêtement $ p : \tilde{X} \to X $ est une surjection continue dont chaque
fibre admet un voisinage uniformément recouvert --- un voisinage dont la préimage est une réunion
disjointe de copies de lui.
\begin{enumerate}
\item[(a)] Démontrez les propriétés de relèvement des chemins et des homotopies pour les
revêtements.
\item[(b)] Déterminez tous les revêtements connexes du cercle, avec démonstration.
\item[(c)] Montrez que $ \mathbb{R}^2 \to T^2 $ (quotient par $ \mathbb{Z}^2 $) et
$ S^2 \to \mathbb{RP}^2 $ (quotient antipodal) sont des revêtements universels --- des revêtements
simplement connexes --- et utilisez le second pour démontrer que
$ \pi_1(\mathbb{RP}^2) \cong \mathbb{Z}/2 $.
\item[(d)] Construisez le revêtement universel du huit, montrez que c'est un arbre $ 4 $-valent infini, et
déduisez-en que $ \pi_1 $ du huit est le groupe libre à deux générateurs.
\item[(e)] Démontrez que pour un revêtement universel, le groupe des automorphismes de revêtement --- les
homéomorphismes de $ \tilde{X} $ commutant avec $ p $ --- est isomorphe à $ \pi_1(X) $.
\end{enumerate}

\end{problem}

\noindent Les revêtements descendent des surfaces de Riemann pour les fonctions multiformes
(1851) et du revêtement universel implicite dans les travaux d'uniformisation de Poincaré ; la
correspondance moderne, à la Galois, entre sous-groupes de $ \pi_1 $ et revêtements a été systématisée
dans les années 1920--30. La question (d) est un petit miracle --- un fait purement algébrique sur les
groupes libres (les sous-groupes d'un groupe libre sont libres, théorème de Nielsen--Schreier) devient
transparent lu à travers la topologie des arbres, un aperçu précoce de la théorie géométrique des
groupes.

\begin{refs}
\item Contexte : Revêtement (mathématiques) --- Wikipédia (\url{https://fr.wikipedia.org/wiki/Rev%C3%AAtement_(math%C3%A9matiques)})
\item Lecture : A. Hatcher, \textit{Algebraic Topology}, ch. 1 --- gratuit sur la page de l'auteur (en anglais) (\url{https://pi.math.cornell.edu/~hatcher/AT/ATpage.html})
\item Lecture : J. R. Munkres, \textit{Topology}, ch. 13
\end{refs}

\bigskip

\begin{problem}[{Van Kampen et les groupes de surfaces --- Master}]
\textbf{TOP-18.} Le théorème de Seifert--Van Kampen calcule $ \pi_1(U \cup V) $ à partir de
$ \pi_1(U) $, $ \pi_1(V) $ et $ \pi_1(U \cap V) $ lorsque $ U, V $ sont ouverts et connexes par
arcs, d'intersection connexe par arcs.
\begin{enumerate}
\item[(a)] Utilisez-le pour montrer que $ \pi_1 $ d'un bouquet de $ n $ cercles est libre à $ n $
générateurs.
\item[(b)] À partir des présentations carrées de TOP-09, calculez les présentations
\[ \pi_1(T^2) = \langle a, b \mid aba^{-1}b^{-1} \rangle, \quad
\pi_1(K) = \langle a, b \mid abab^{-1} \rangle, \quad
\pi_1(\mathbb{RP}^2) = \langle a \mid a^2 \rangle \cong \mathbb{Z}/2 , \]
de sorte qu'en particulier le groupe du tore est $ \mathbb{Z}^2 $.
\item[(c)] Calculez les abélianisés des groupes fondamentaux de toutes les surfaces compactes (orientables de
genre $ g $, non orientables à $ k $ bonnets croisés) et concluez qu'aucune paire de surfaces de la
liste de classification de TOP-19 n'est formée de surfaces homéomorphes --- achevant ainsi la moitié
\og{} distinction \fg{} de ce théorème.
\end{enumerate}

\end{problem}

\noindent Seifert (1931) et van Kampen (1933) ont démontré indépendamment des versions du
théorème de recollement ; c'est le principal moteur de calcul du groupe fondamental, convertissant
directement les présentations par découpage-collage des surfaces en générateurs et relations. La question
(c) boucle une boucle : les invariants abélianisés sont exactement les premiers groupes d'homologie de
Poincaré, si bien que le problème redécouvre discrètement l'homologie une page avant que TOP-20 ne la
définisse comme il faut. Les groupes de surfaces eux-mêmes ont connu une seconde vie comme objets
fondamentaux de la théorie géométrique des groupes.

\begin{refs}
\item Contexte : Théorème de Seifert--Van Kampen --- Wikipédia (\url{https://fr.wikipedia.org/wiki/Th%C3%A9or%C3%A8me_de_van_Kampen})
\item Lecture : A. Hatcher, \textit{Algebraic Topology}, ch. 1 --- gratuit sur la page de l'auteur (en anglais) (\url{https://pi.math.cornell.edu/~hatcher/AT/ATpage.html})
\item Lecture : J. R. Munkres, \textit{Topology}, ch. 11
\end{refs}

\bigskip

\begin{problem}[{La classification des surfaces compactes --- Master}]
\textbf{TOP-19.} Démontrez le théorème de classification : toute surface compacte connexe (variété
de dimension 2 sans bord, supposée triangulable) est homéomorphe à exactement l'une des surfaces
suivantes --- la sphère, une somme connexe de $ g \ge 1 $ tores, ou une somme connexe de $ k \ge 1 $
plans projectifs. Suivez la voie classique :
\begin{enumerate}
\item[(a)] montrez que toute surface compacte triangulée admet une présentation polygonale (un seul polygone
dont les arêtes sont identifiées deux à deux) ;
\item[(b)] ramenez une telle présentation à une forme normale par des mouvements de découpage-collage ;
\item[(c)] vérifiez les caractéristiques d'Euler $ \chi = 2 - 2g $ et $ \chi = 2 - k $, et démontrez que
la caractéristique d'Euler jointe à l'orientabilité distingue toutes les formes normales (TOP-18(c)
fournit la moitié théorie des groupes).
\end{enumerate}

\end{problem}

\noindent La classification est la première réponse complète de la topologie à une question du
type \og{} dressez la liste de tous les espaces \fg{}. Möbius a classifié les surfaces orientables en 1863 ; Dyck
et d'autres ont traité le cas non orientable ; des démonstrations pleinement rigoureuses sont venues avec
Dehn et Heegaard (1907) et avec le théorème de Radó de 1925 selon lequel toute surface est triangulable.
La complétude à deux invariants du théorème --- genre et orientabilité suffisent --- a fixé le modèle de toute
classification rêvée depuis, et l'absence de toute liste comparable en dimension trois est exactement ce
qui a rendu la conjecture de Poincaré (voir la bande Recherche) si difficile.

\begin{refs}
\item Contexte : Classification des surfaces --- Wikipédia (\url{https://fr.wikipedia.org/wiki/Classification_des_surfaces})
\item Contexte : Genre (mathématiques) --- Wikipédia (\url{https://fr.wikipedia.org/wiki/Genre_(math%C3%A9matiques)})
\item Lecture : M. A. Armstrong, \textit{Basic Topology}, ch. 7
\item Lecture : J. R. Munkres, \textit{Topology}, ch. 12
\end{refs}

\bigskip

\begin{problem}[{Premiers pas en homologie simpliciale --- Master}]
\textbf{TOP-20.} Pour un complexe simplicial, le $ n $-ième groupe de chaînes est le groupe
abélien libre sur les $ n $-simplexes ; les applications de bord $ \partial_n $ envoient un simplexe
sur la somme alternée de ses faces ; l'homologie est
\[ H_n = \ker \partial_n / \operatorname{im} \partial_{n+1} . \]
\begin{enumerate}
\item[(a)] Démontrez que $ \partial_{n} \circ \partial_{n+1} = 0 $, de sorte que la définition ait un
sens.
\item[(b)] Calculez tous les groupes d'homologie du cercle, de la sphère $ S^2 $, du tore, de la bouteille
de Klein et de $ \mathbb{RP}^2 $ à partir de triangulations explicites (ou de structures de
$\Delta$-complexes), et confrontez les réponses à l'intuition : $ H_0 $ compte les composantes, $ H_1 $
compte les lacets indépendants, $ H_2 $ détecte le volume enclos et l'orientabilité.
\item[(c)] Démontrez la formule d'Euler--Poincaré : la somme alternée des rangs des groupes d'homologie vaut
$ V - E + F $ pour toute triangulation --- la caractéristique d'Euler de TOP-02 est donc indépendante
de la triangulation.
\item[(d)] Expliquez comment l'homologie en dimension $ n - 1 $ démontre le théorème de non-rétraction pour
$ D^n $, achevant TOP-14(d).
\end{enumerate}

\end{problem}

\noindent Poincaré a introduit la machinerie (sous forme de nombres de Betti et de coefficients
de torsion) dans \textit{Analysis Situs} ; Emmy Noether a observé dans les années 1920 que ces nombres
étaient les ombres de groupes abéliens, et la forme moderne est due à son entourage. L'homologie est plus
difficile à définir que le groupe fondamental mais bien plus facile à calculer --- un échange que démontre la
question (b) --- et elle fonctionne en toute dimension, ce qu'exploite la question (d). Cette porte ouvre
sur l'homologie singulière, la cohomologie et la topologie algébrique des chapitres 2 et 3 de
Hatcher.

\begin{refs}
\item Contexte : Homologie simpliciale --- Wikipédia (en anglais) (\url{https://en.wikipedia.org/wiki/Simplicial_homology})
\item Contexte : Complexe simplicial --- Wikipédia (\url{https://fr.wikipedia.org/wiki/Complexe_simplicial}), Homologie (mathématiques) --- Wikipédia (\url{https://fr.wikipedia.org/wiki/Homologie_(math%C3%A9matiques)})
\item Lecture : A. Hatcher, \textit{Algebraic Topology}, ch. 2 --- gratuit sur la page de l'auteur (en anglais) (\url{https://pi.math.cornell.edu/~hatcher/AT/ATpage.html})
\item Cours : MIT OCW 18.905 Algebraic Topology I (en anglais) (\url{https://ocw.mit.edu/courses/18-905-algebraic-topology-i-fall-2016/})
\end{refs}

\bigskip

\begin{problem}[{Le n\oe{}ud de trèfle est vraiment noué --- Master}]
\textbf{TOP-21.} Un n\oe{}ud est un cercle plongé dans $ \mathbb{R}^3 $, étudié à travers ses
diagrammes ; deux diagrammes représentent le même n\oe{}ud exactement lorsqu'ils sont reliés par des
déformations planes et les trois mouvements de Reidemeister. On dit qu'un diagramme est
\textit{tricoloriable} si chaque arc peut être colorié avec l'une de trois couleurs de sorte qu'à chaque
croisement les trois arcs qui se rencontrent aient toutes les couleurs identiques ou toutes différentes,
et qu'au moins deux couleurs soient employées.
\begin{enumerate}
\item[(a)] Démontrez que la tricoloriabilité est préservée par chacun des trois mouvements de Reidemeister ---
c'est un invariant de n\oe{}ud.
\item[(b)] Montrez que le diagramme standard du n\oe{}ud de trèfle est tricoloriable et que le diagramme rond du
n\oe{}ud trivial ne l'est pas, et concluez que le trèfle ne peut pas être dénoué.
\item[(c)] Montrez que le n\oe{}ud de huit n'est \textit{pas} tricoloriable, et expliquez ce que cela vous apprend
et ce que cela ne vous apprend pas.
\item[(d)] Généralisez : définissez la $ p $-coloriabilité pour un nombre premier impair $ p $, démontrez
l'invariance, et trouvez un $ p $ qui distingue le n\oe{}ud de huit du n\oe{}ud trivial.
\end{enumerate}

\end{problem}

\noindent La théorie des n\oe{}uds a commencé avec la spéculation de Kelvin en 1867 selon laquelle
les atomes seraient des tubes de tourbillon noués, ce qui a lancé Tait dans la tabulation des n\oe{}uds ; les
mathématiques ont survécu à la physique. Reidemeister a démontré en 1927 que ses mouvements suffisent,
rendant l'équivalence de n\oe{}uds combinatoire, et l'astuce du tricoloriage --- avatar des représentations du
groupe du n\oe{}ud, introduit par Ralph Fox vers 1956 --- est la démonstration honnête la plus simple qu'un n\oe{}ud
est non trivial. La question (d) ouvre la porte à l'industrie moderne des invariants de n\oe{}uds : Alexander,
Jones et au-delà (voir TOP-23 et TOP-24).

\begin{refs}
\item Contexte : Tricolorabilité --- Wikipédia (\url{https://fr.wikipedia.org/wiki/Tricolorabilit%C3%A9})
\item Contexte : N\oe{}ud de trèfle --- Wikipédia (\url{https://fr.wikipedia.org/wiki/N%C5%93ud_de_tr%C3%A8fle}), Mouvements de Reidemeister --- Wikipédia (\url{https://fr.wikipedia.org/wiki/Mouvements_de_Reidemeister})
\item Lecture : C. C. Adams, \textit{The Knot Book}, ch. 1
\end{refs}

\bigskip

\begin{problem}[{Ajouter un point à l'infini --- Master, short}]
\textbf{TOP-31.} Pour un espace localement compact séparé $ X $, soit
$ X^* = X \cup \{\infty\} $ muni de la topologie dont les ouverts sont les ouverts de $ X $ ainsi
que les ensembles $ X^* \setminus K $ pour $ K \subseteq X $ compact et fermé. Démontrez que
$ X^* $ est compact séparé et que $ (\mathbb{R}^n)^* $ est homéomorphe à la sphère
$ S^n $.
\end{problem}

\noindent Pavel Alexandroff a introduit cette construction en 1924 : un point de plus, et un
espace qui filait à l'infini devient compact. L'homéomorphisme de la seconde moitié est la projection
stéréographique lue à l'envers, et c'est la raison pour laquelle $ \mathbb{C} \cup \{\infty\} $
s'appelle la sphère de Riemann --- l'image la plus utile de toute l'analyse complexe.

\begin{refs}
\item Contexte : Compactifié d'Alexandrov --- Wikipédia (\url{https://fr.wikipedia.org/wiki/Compactifi%C3%A9_d'Alexandrov})
\item Lecture : J. R. Munkres, \textit{Topology}, \S{}29
\end{refs}

\bigskip

\begin{problem}[{Une application qui rate un point est de degré zéro --- Master, short}]
\textbf{TOP-32.} Soit $ f : S^n \to S^n $ continue et non surjective. Démontrez
que $ f $ est homotope à une application constante, et concluez que $ \deg f = 0 $.
\end{problem}

\noindent Le degré d'une application de la sphère dans elle-même --- l'entier par lequel $ f $
multiplie la classe d'homologie de tête --- est l'invariant de Brouwer, datant des environs de 1910, et ce
lemme de deux lignes est le levier qui le transforme en géométrie. C'est l'étape qui fait fonctionner les
démonstrations standard du théorème de la boule chevelue (TOP-15) et du théorème du point fixe de Brouwer
(TOP-14) : une application sans point fixe rate quelque chose, et une application qui rate quelque chose ne
peut pas être de degré $ (-1)^{n+1} $.

\begin{refs}
\item Contexte : Degré d'une application continue --- Wikipédia (\url{https://fr.wikipedia.org/wiki/Degr%C3%A9_d'une_application})
\item Lecture : A. Hatcher, \textit{Algebraic Topology}, \S{}2.2 --- gratuit sur la page de l'auteur (en anglais) (\url{https://pi.math.cornell.edu/~hatcher/AT/ATpage.html})
\end{refs}

\bigskip

\begin{problem}[{Le lemme d'Urysohn, et ce qu'achète la séparation --- Master}]
\textbf{TOP-33.} On dit qu'un espace topologique $ X $ est \textit{normal} si les singletons sont
fermés et si deux fermés disjoints quelconques admettent des voisinages ouverts disjoints. Rien dans
cette hypothèse ne mentionne les nombres réels.
\begin{enumerate}
\item[(a)] Démontrez le lemme d'Urysohn : si $ A $ et $ B $ sont des fermés disjoints d'un espace normal
$ X $, il existe une application continue $ f : X \to [0,1] $ avec $ f \equiv 0 $ sur $ A $ et
$ f \equiv 1 $ sur $ B $. (Construisez une famille d'ouverts indexée par les rationnels dyadiques de
$ [0,1] $, emboîtés de sorte que les adhérences respectent l'ordre, et lisez $ f $ sur la
famille.)
\item[(b)] Montrez que dans un espace métrique la fonction
\[ f(x) = \frac{d(x,A)}{d(x,A) + d(x,B)} \]
fait immédiatement l'affaire, et expliquez ce qui manque à un espace normal général et rend la question
(a) tellement plus longue.
\item[(c)] Déduisez-en le théorème de prolongement de Tietze : une fonction continue à valeurs réelles sur un
sous-espace fermé d'un espace normal se prolonge en une fonction continue sur l'espace entier.
\item[(d)] Démontrez le théorème de métrisation d'Urysohn : tout espace régulier à base dénombrable se plonge
dans le cube de Hilbert $ [0,1]^{\mathbb{N}} $, et est donc métrisable --- bouclant la boucle ouverte
par TOP-05(d).
\end{enumerate}

\end{problem}

\noindent Pavel Urysohn s'est noyé en 1924 en nageant au large de la Bretagne, à 26 ans ; le
lemme a paru à titre posthume, édité par Alexandroff, et c'est le moment où la topologie générale cesse
d'être de la comptabilité. Munkres l'appelle le premier théorème profond de son livre, et la raison est la
question (b) : les axiomes de séparation sont de pures hypothèses ensemblistes sur les ouverts, et pourtant
ils fabriquent une fonction continue à valeurs réelles à partir de rien. Tout ce qui permet à la topologie
de parler à l'analyse --- partitions de l'unité, plongements dans des cubes, théorèmes de prolongement ---
passe par cette construction.

\begin{refs}
\item Contexte : Lemme d'Urysohn --- Wikipédia (\url{https://fr.wikipedia.org/wiki/Lemme_d'Urysohn})
\item Contexte : Théorème de prolongement de Tietze --- Wikipédia (\url{https://fr.wikipedia.org/wiki/Th%C3%A9or%C3%A8me_de_prolongement_de_Tietze}), Espace métrisable (théorème de métrisation d'Urysohn) --- Wikipédia (\url{https://fr.wikipedia.org/wiki/Espace_m%C3%A9trisable})
\item Lecture : J. R. Munkres, \textit{Topology}, \S{}33--34, \S{}40
\item Cours : MIT OCW 18.901 Introduction to Topology (en anglais) (\url{https://ocw.mit.edu/courses/18-901-introduction-to-topology-fall-2004/})
\end{refs}

\bigskip

\begin{problem}[{L'application de Hopf est essentielle --- Master, short}]
\textbf{TOP-42.} Soit $ h : S^3 \to S^2 $ l'application de Hopf, définie sur les
vecteurs unitaires $ (z_1, z_2) \in \mathbb{C}^2 $ par
\[ h(z_1, z_2) = [\, z_1 : z_2 \,] \in \mathbb{CP}^1 \cong S^2 . \]
Démontrez que $ h $ n'est pas homotope à une application constante, et concluez que
$ \pi_3(S^2) \neq 0 $.
\end{problem}

\noindent Heinz Hopf a publié cette application en 1931, et elle a renversé l'attente régnante :
l'homologie d'une sphère s'annule au-dessus de sa propre dimension, si bien que tout le monde supposait
qu'il en irait de même pour l'homotopie. Il n'en est rien. Chaque fibre $ h^{-1}(p) $ est un grand
cercle, deux fibres quelconques sont enlacées exactement une fois, et ce nombre d'enlacement ---
l'invariant de Hopf --- est ce qui obstrue l'homotopie nulle ; la même application présente $ S^3 $ comme
un fibré de fibre $ S^1 $ au-dessus de $ S^2 $. La réponse complète est
$ \pi_3(S^2) \cong \mathbb{Z} $, engendré par $ h $, et le calcul de $ \pi_{n+k}(S^n) $ en général
reste l'un des grands chantiers inachevés de la topologie.

\begin{refs}
\item Contexte : Fibration de Hopf --- Wikipédia (\url{https://fr.wikipedia.org/wiki/Fibration_de_Hopf})
\item Contexte : Groupes d'homotopie des sphères --- Wikipédia (\url{https://fr.wikipedia.org/wiki/Groupes_d'homotopie_des_sph%C3%A8res}), Fibré --- Wikipédia (\url{https://fr.wikipedia.org/wiki/Fibr%C3%A9})
\item Lecture : A. Hatcher, \textit{Algebraic Topology}, \S{}4.2 --- gratuit sur la page de l'auteur (en anglais) (\url{https://pi.math.cornell.edu/~hatcher/AT/ATpage.html})
\end{refs}

\bigskip

\begin{problem}[{Couper les espaces en deux : Mayer--Vietoris --- Master}]
\textbf{TOP-43.} Soit $ X = A \cup B $ avec $ A, B $ ouverts (ou une paire CW avec les
ajustements évidents). La suite de Mayer--Vietoris est la suite exacte longue
\[ \cdots \to H_n(A \cap B) \to H_n(A) \oplus H_n(B) \to H_n(X)
\xrightarrow{\ \partial\ } H_{n-1}(A \cap B) \to \cdots \]
\begin{enumerate}
\item[(a)] Établissez-la : écrivez la suite exacte courte de complexes de chaînes dont elle provient,
expliquez pourquoi le terme du milieu peut être remplacé par des chaînes petites relativement au
recouvrement (c'est là qu'interviennent la subdivision barycentrique et le lemme du nombre de Lebesgue
de TOP-29), et appliquez le lemme du serpent.
\item[(b)] Calculez $ H_*(S^n) $ pour tout $ n $ par récurrence, en recouvrant la sphère par deux ouverts
contractiles dont l'intersection se rétracte par déformation sur $ S^{n-1} $.
\item[(c)] Calculez l'homologie de la surface fermée orientable $ \Sigma_g $ de genre $ g $ en la coupant
en un disque et un voisinage d'un bouquet de $ 2g $ cercles, obtenant les nombres de Betti
$ 1, 2g, 1 $. Faites de même pour la surface non orientable à $ k $ bonnets croisés et trouvez où
apparaît la torsion.
\item[(d)] Soit $ K \subseteq S^3 $ un n\oe{}ud de voisinage tubulaire $ N $. Appliquez la suite à
$ A = N $ et $ B = S^3 \setminus K $ --- leur intersection est un tore --- et démontrez que
$ H_1(S^3 \setminus K) \cong \mathbb{Z} $ pour \textit{tout} n\oe{}ud. Concluez que l'homologie seule ne peut
jamais distinguer le trèfle du n\oe{}ud trivial, et dites ce que TOP-21 utilise à la place.
\end{enumerate}

\end{problem}

\noindent Walther Mayer a traité un cas particulier en 1929 et Leopold Vietoris le cas général
peu après --- Vietoris, un Autrichien qui publiait encore à plus de quatre-vingt-dix ans, a vécu jusqu'à
110 ans. Leur suite est le cheval de trait calculatoire de l'homologie : elle transforme le geste
géométrique consistant à couper un espace en deux morceaux en une recette algébrique, exactement comme
Seifert--van Kampen (TOP-18) le fait pour le groupe fondamental. La question (d) est le cas le plus simple
de la dualité d'Alexander, et sa conclusion négative est la raison pour laquelle la théorie des n\oe{}uds a dû
inventer des invariants plus fins que l'homologie.

\begin{refs}
\item Contexte : Suite de Mayer--Vietoris --- Wikipédia (\url{https://fr.wikipedia.org/wiki/Suite_de_Mayer-Vietoris})
\item Contexte : Nombre de Betti --- Wikipédia (\url{https://fr.wikipedia.org/wiki/Nombre_de_Betti}), Dualité d'Alexander --- Wikipédia (\url{https://fr.wikipedia.org/wiki/Dualit%C3%A9_d'Alexander})
\item Lecture : A. Hatcher, \textit{Algebraic Topology}, \S{}2.2 --- gratuit sur la page de l'auteur (en anglais) (\url{https://pi.math.cornell.edu/~hatcher/AT/ATpage.html})
\item Cours : MIT OCW 18.905 Algebraic Topology I (en anglais) (\url{https://ocw.mit.edu/courses/18-905-algebraic-topology-i-fall-2016/})
\end{refs}

\bigskip

\begin{problem}[{Compter les points fixes avec le nombre de Lefschetz --- Master}]
\textbf{TOP-44.} Pour un CW-complexe fini $ X $ (TOP-41) et une application continue
$ f : X \to X $, on définit le nombre de Lefschetz
\[ \Lambda(f) = \sum_{n \ge 0} (-1)^n \operatorname{tr}\Big( f_* : H_n(X;\mathbb{Q}) \to H_n(X;\mathbb{Q}) \Big) . \]
Le théorème du point fixe de Lefschetz affirme que si $ \Lambda(f) \neq 0 $ alors $ f $ a un point
fixe. On admet le théorème à partir de la question (b).
\begin{enumerate}
\item[(a)] Démontrez que $ \Lambda(\mathrm{id}_X) = \chi(X) $, et déduisez-en que sur un espace de
caractéristique d'Euler non nulle toute application homotope à l'identité a un point fixe.
\item[(b)] Démontrez qu'un CW-complexe fini dont l'homologie rationnelle est celle d'un point a la propriété
du point fixe, et déduisez-en le théorème du point fixe de Brouwer pour $ D^n $ --- une seconde voie
vers TOP-14(d).
\item[(c)] Pour $ f : S^n \to S^n $ de degré $ d $, calculez $ \Lambda(f) $. Déduisez-en qu'une
application de $ S^n $ dans elle-même sans point fixe est de degré $ (-1)^{n+1} $, et utilisez cela
pour démontrer le théorème de la boule chevelue (TOP-15) pour toute sphère de dimension paire.
\item[(d)] Démontrez que toute application continue de $ \mathbb{RP}^{2n} $ dans lui-même a un point fixe,
et montrez que cela échoue en dimension impaire en exhibant une application de
$ \mathbb{RP}^{2n+1} $ dans lui-même sans point fixe.
\item[(e)] Montrez que le théorème ne donne aucune information dans l'autre sens : exhibez une application
avec $ \Lambda(f) = 0 $ qui a des points fixes, et une avec $ \Lambda(f) = 0 $ qui n'en a
pas.
\end{enumerate}

\end{problem}

\noindent Solomon Lefschetz a démontré sa formule en 1926, et Hopf l'a bientôt étendue à tous les
complexes finis ; c'est le théorème qui a rendu calculables les questions de points fixes, puisque
$ \Lambda $ ne dépend que des applications induites en homologie, lesquelles ne dépendent que de la
classe d'homotopie de $ f $. Le théorème de Brouwer est le cas particulier où l'espace n'a pour ainsi
dire pas d'homologie. L'idée ne s'est pas arrêtée là : le théorème du point fixe d'Atiyah et Bott (1966)
est le même énoncé pour les complexes elliptiques, et la démonstration des conjectures de Weil par
Grothendieck compte les points d'une variété sur un corps fini comme les points fixes de Lefschetz du
morphisme de Frobenius.

\begin{refs}
\item Contexte : Théorème du point fixe de Lefschetz --- Wikipédia (\url{https://fr.wikipedia.org/wiki/Th%C3%A9or%C3%A8me_du_point_fixe_de_Lefschetz})
\item Contexte : Degré d'une application continue --- Wikipédia (\url{https://fr.wikipedia.org/wiki/Degr%C3%A9_d'une_application})
\item Lecture : A. Hatcher, \textit{Algebraic Topology}, \S{}2.C --- gratuit sur la page de l'auteur (en anglais) (\url{https://pi.math.cornell.edu/~hatcher/AT/ATpage.html})
\item Cours : MIT OCW 18.905 Algebraic Topology I (en anglais) (\url{https://ocw.mit.edu/courses/18-905-algebraic-topology-i-fall-2016/})
\end{refs}

\bigskip


\section*{Recherche}

\begin{problem}[{La conjecture de Poincaré lisse en dimension quatre --- Recherche}]
\textbf{TOP-22.} (Ouvert.) La question de Poincaré lisse dans la seule dimension où elle tient
encore.
\begin{enumerate}
\item[(a)] Toute variété lisse de dimension $ 4 $ homotopiquement équivalente à la sphère $ 4 $-dimensionnelle
est-elle réellement \textit{difféomorphe} à $ S^4 $ ? De manière équivalente : $ S^4 $ admet-elle une
structure lisse exotique --- une variété lisse qui lui est homéomorphe mais non difféomorphe ?
\item[(b)] Une tâche d'entrée abordable : rédigez un panorama soigné expliquant pourquoi la dimension quatre
est différente, en exposant ce que l'on sait dans toutes les autres dimensions et où les outils standard
de dimension supérieure (théorème du h-cobordisme, astuce de Whitney) s'effondrent en dimension
quatre.
\end{enumerate}

\end{problem}

\noindent La conjecture de Poincaré topologique est réglée en toute dimension : Smale a démontré
les dimensions $ \ge 5 $ (1961), Freedman la dimension quatre (1982) et Perelman la dimension trois
(2003, via le flot de Ricci --- le cas du prix du millénaire, traité sur notre page Problèmes ouverts). La
question \textit{lisse} en dimension quatre est la dernière debout, et en 2026 elle est ouverte. La dimension
quatre est réellement étrange : les travaux de Freedman et de Donaldson ont ensemble produit des structures
lisses exotiques sur $ \mathbb{R}^4 $ --- sur aucun autre $ \mathbb{R}^n $ il n'en existe --- et de
nombreuses candidates $ 4 $-sphères exotiques (sphères de Cappell--Shaneson, torsions de Gluck) ont été
proposées, de grandes familles s'avérant ensuite standard (Akbulut 2010, Gompf). Les approches par la
théorie des n\oe{}uds (Freedman--Gompf--Morrison--Walker 2010 ; Manolescu--Piccirillo 2021) restent actives ; des
démonstrations périodiquement annoncées n'ont jusqu'ici pas résisté à l'examen, et il n'existe ni
démonstration consensuelle ni contre-exemple.

\begin{refs}
\item Contexte : Conjecture de Poincaré généralisée --- Wikipédia (en anglais) (\url{https://en.wikipedia.org/wiki/Generalized_Poincar%C3%A9_conjecture})
\item Contexte : Variété de dimension 4 --- Wikipédia (en anglais) (\url{https://en.wikipedia.org/wiki/4-manifold}), R4 exotique --- Wikipédia (en anglais) (\url{https://en.wikipedia.org/wiki/Exotic_R4})
\item Article : M. Freedman, \og{} The topology of four-dimensional manifolds \fg{} (1982), \textit{Journal of Differential Geometry} 17
\end{refs}

\bigskip

\begin{problem}[{La conjecture bordant-ruban --- Recherche}]
\textbf{TOP-23.} (Ouvert.) Un n\oe{}ud de $ S^3 = \partial B^4 $ est \textit{bordant} (slice) s'il borde
un disque lissement plongé dans la boule de dimension quatre $ B^4 $, et \textit{ruban} s'il borde un
disque immergé dans $ S^3 $ dont les seules auto-intersections sont des singularités en ruban --- des
arcs où une nappe traverse proprement une autre. La conjecture bordant-ruban de Fox affirme que ces deux
notions coïncident.
\begin{enumerate}
\item[(a)] Démontrez que tout n\oe{}ud ruban est bordant.
\item[(b)] Montrez que la somme connexe d'un n\oe{}ud quelconque avec son image miroir inversée est ruban (les
n\oe{}uds bordants sont donc abondants).
\item[(c)] Étudiez pourquoi le caractère bordant est véritablement quadridimensionnel.
\item[(d)] La conjecture elle-même : démontrez que tout n\oe{}ud bordant est ruban, ou trouvez un
contre-exemple.
\end{enumerate}

\end{problem}

\noindent Ralph Fox a posé la question en 1962, et elle est devenue un problème organisateur
central de la théorie des n\oe{}uds en dimension quatre, situé juste sous TOP-22 --- le caractère bordant de
n\oe{}uds particuliers est intimement lié à la recherche de sphères de dimension quatre exotiques. La
conjecture est démontrée pour de grandes familles : n\oe{}uds à deux ponts (Lisca 2007, via une théorie de
jauge à la Donaldson), n\oe{}uds pretzel à trois brins (Greene--Jabuka 2011), et d'autres ; des
contre-exemples potentiels sont régulièrement proposés puis écartés. En 2026 la conjecture générale est
ouverte dans les deux directions --- ni démonstration, ni contre-exemple.

\begin{refs}
\item Contexte : N\oe{}ud bordant --- Wikipédia (\url{https://fr.wikipedia.org/wiki/N%C5%93ud_bordant}), N\oe{}ud ruban --- Wikipédia (en anglais) (\url{https://en.wikipedia.org/wiki/Ribbon_knot})
\item Article : P. Lisca, \og{} Lens spaces, rational balls and the ribbon conjecture \fg{} (2007), \textit{Geometry \& Topology} 11
\item Article : J. Greene et S. Jabuka, \og{} The slice-ribbon conjecture for 3-stranded pretzel knots \fg{} (2011), arXiv:0706.3398 (\url{https://arxiv.org/abs/0706.3398})
\item Lecture : C. C. Adams, \textit{The Knot Book}
\end{refs}

\bigskip

\begin{problem}[{Reconnaître le n\oe{}ud trivial : à quel prix ? --- Recherche}]
\textbf{TOP-24.} (Partiellement résolu ; la question centrale est ouverte.) RECONNAISSANCE DU N\OE{}UD
TRIVIAL est le problème de décision suivant : étant donné un diagramme de n\oe{}ud à $ n $ croisements,
est-ce un diagramme du n\oe{}ud trivial ?
\begin{enumerate}
\item[(a)] Montrez que le problème est décidable en principe à partir des mouvements de Reidemeister
\textit{si l'on dispose} d'une borne supérieure sur le nombre de mouvements nécessaires, et expliquez
pourquoi aucune telle borne n'est évidente.
\item[(b)] Parcourez la raison pour laquelle la théorie des surfaces normales de Haken fournit un
algorithme.
\item[(c)] La question centrale : déterminez la complexité algorithmique du problème --- en particulier, s'il
existe un algorithme en temps polynomial.
\end{enumerate}

\end{problem}

\noindent Haken a donné le premier algorithme en 1961 via les surfaces normales. Hass, Lagarias
et Pippenger ont démontré en 1999 que le problème est dans NP ; Kuperberg (2014) a montré l'appartenance à
co-NP en supposant l'hypothèse de Riemann généralisée, et Lackenby (2016, publié en 2021) a retiré cette
hypothèse, de sorte que la trivialité admet des certificats courts dans les deux directions --- indice fort,
au vu des croyances standard en complexité, que le problème n'est pas NP-complet. En 2021 Lackenby a
annoncé un algorithme de reconnaissance du n\oe{}ud trivial s'exécutant en temps quasi polynomial
$ n^{O(\log n)} $. En 2026, aucun algorithme en temps polynomial n'est connu et aucun n'est démontré
impossible ; savoir si RECONNAISSANCE DU N\OE{}UD TRIVIAL est dans P est un problème authentiquement ouvert, à
cheval sur la topologie et l'informatique.

\begin{refs}
\item Contexte : Problème de reconnaissance du n\oe{}ud trivial --- Wikipédia (en anglais) (\url{https://en.wikipedia.org/wiki/Unknotting_problem})
\item Article : J. Hass, J. C. Lagarias, N. Pippenger, \og{} The computational complexity of knot and link problems \fg{} (1999), \textit{Journal of the ACM} 46
\item Article : M. Lackenby, \og{} The efficient certification of knottedness and Thurston norm \fg{} (2021), \textit{Advances in Mathematics}, arXiv:1604.00290 (\url{https://arxiv.org/abs/1604.00290})
\item Article : M. Lackenby, \og{} Unknot recognition in quasi-polynomial time \fg{} (annoncé en 2021), diapositives d'exposé, Oxford (en anglais) (\url{https://people.maths.ox.ac.uk/lackenby/quasipolynomial-talk-oxford-compressed.pdf})
\end{refs}

\bigskip

\begin{problem}[{La question d'asphéricité de Whitehead --- Recherche, short}]
\textbf{TOP-34.} (Ouvert.) On dit qu'un CW-complexe connexe est \textit{asphérique} si
$ \pi_n = 0 $ pour tout $ n \ge 2 $, autrement dit si son revêtement universel est contractile.
Tranchez la question posée par J. H. C. Whitehead en 1941 : tout sous-complexe connexe d'un CW-complexe
asphérique de dimension deux est-il lui-même asphérique ?
\end{problem}

\noindent Whitehead a soulevé la question à la fin d'un article sur l'ajout de relations aux
groupes d'homotopie, et elle résiste depuis --- un rare problème ouvert dont l'énoncé n'exige rien de plus
qu'un premier cours de topologie algébrique. James Howie l'a réduite dans les années 1980 à deux
configurations particulières, et elle est connue pour de larges classes de complexes (par exemple ceux dont
le groupe fondamental est de présentation finie et localement indicable), mais en 2026 ni démonstration ni
contre-exemple n'est connu ; un contre-exemple réglerait aussi des questions voisines en théorie
combinatoire des groupes.

\begin{refs}
\item Contexte : Conjecture de Whitehead --- Wikipédia (en anglais) (\url{https://en.wikipedia.org/wiki/Whitehead_conjecture}), Espace asphérique --- Wikipédia (\url{https://fr.wikipedia.org/wiki/Espace_asph%C3%A9rique})
\item Article : J. H. C. Whitehead, \og{} On adding relations to homotopy groups \fg{} (1941), \textit{Annals of Mathematics} 42
\end{refs}

\bigskip

\begin{problem}[{La conjecture d'effondrement de Zeeman --- Recherche, short}]
\textbf{TOP-35.} (Ouvert.) Un polyèdre fini \textit{s'effondre} s'il peut être réduit
à un point en supprimant à répétition une face libre avec l'unique cellule qui la contient. Tranchez la
conjecture de Zeeman : pour tout polyèdre contractile $ K $ de dimension deux, le produit
$ K \times I $ est-il effondrable, où $ I = [0,1] $ ?
\end{problem}

\noindent Contractile n'implique pas effondrable --- le bonnet d'âne de Zeeman lui-même (1964), un
triangle dont les trois arêtes sont recollées selon un cycle mal appareillé, en est le témoin standard, et
pourtant le bonnet d'âne multiplié par un intervalle, lui, s'effondre. Zeeman a conjecturé qu'il en va
toujours ainsi ; sa conjecture implique à la fois la conjecture de Poincaré (réglée par Perelman en 2003)
et la conjecture d'Andrews--Curtis (toujours ouverte), ce qui donne une juste mesure de sa difficulté. Des
résultats partiels sont connus --- le produit d'un 2-polyèdre contractile par un cube d'assez grande
dimension s'effondre --- mais en 2026 la conjecture elle-même est ouverte.

\begin{refs}
\item Contexte : Conjecture de Zeeman --- Wikipédia (en anglais) (\url{https://en.wikipedia.org/wiki/Zeeman_conjecture}), Bonnet d'âne (topologie) --- Wikipédia (en anglais) (\url{https://en.wikipedia.org/wiki/Dunce_hat_(topology)})
\item Article : E. C. Zeeman, \og{} On the dunce hat \fg{} (1964), \textit{Topology} 2
\end{refs}

\bigskip

\begin{problem}[{La conjecture de Novikov --- Recherche}]
\textbf{TOP-36.} (Ouvert.) Soit $ M $ une variété lisse close orientée de groupe fondamental
$ \pi $, soit $ L(M) $ sa $ L $-classe de Hirzebruch, et soit $ u : M \to B\pi $ l'application
classifiant le revêtement universel. Les \textit{signatures supérieures} de $ M $ sont les nombres
rationnels
\[ \sigma_x(M) \;=\; \big\langle\, L(M) \smile u^*(x),\; [M] \,\big\rangle , \qquad
x \in H^*(B\pi; \mathbb{Q}) . \]
\begin{enumerate}
\item[(a)] Prenez $ \dim M = 4k $, de sorte que la signature de la forme d'intersection sur
$ H^{2k}(M;\mathbb{R}) $ soit définie. Vérifiez que pour $ \pi $ trivial la seule signature
supérieure est $ \langle L(M), [M] \rangle $, et énoncez le théorème de la signature de Hirzebruch,
selon lequel ce nombre est exactement cette signature.
\item[(b)] Démontrez que la signature elle-même est un invariant d'équivalence d'homotopie orientée, en
utilisant seulement qu'elle est déterminée par la forme cup-produit sur la cohomologie de dimension
médiane.
\item[(c)] Expliquez pourquoi la conjecture n'est pas une formalité : les classes de Pontryagin rationnelles,
et donc $ L(M) $, ne sont \textit{pas} des invariants d'homotopie en général, bien que Novikov ait
démontré en 1965 que ce sont des invariants topologiques.
\item[(d)] La conjecture : démontrez que toute signature supérieure est un invariant d'équivalence d'homotopie
orientée, pour tout groupe $ \pi $ --- ou exhibez un groupe et une équivalence d'homotopie pour
lesquels un $ \sigma_x $ change.
\end{enumerate}

\end{problem}

\noindent Sergueï Novikov a formulé la conjecture au milieu des années 1960, et elle est devenue
le point de rencontre de la théorie de la chirurgie, de la théorie géométrique des groupes et des algèbres
d'opérateurs : les attaques standard passent par l'application d'assemblage du tableau de Baum--Connes
plutôt que par les variétés. C'est un théorème pour d'énormes classes de groupes --- Lusztig pour les groupes
abéliens libres, Mishchenko et Kasparov pour les sous-groupes discrets de groupes de Lie, Connes et
Moscovici pour les groupes hyperboliques au sens des mots (1990), Higson--Kasparov pour les groupes ayant la
propriété de Haagerup, et Guoliang Yu (2000) pour tout groupe qui se plonge grossièrement dans un espace de
Hilbert. Les groupes aléatoires de Gromov ont montré que tous les groupes ne se plongent pas ainsi, si bien
que la voie de Yu n'achève pas le travail. En 2026 la conjecture est ouverte en général et aucun
contre-exemple n'est connu.

\begin{refs}
\item Contexte : Conjecture de Novikov --- Wikipédia (en anglais) (\url{https://en.wikipedia.org/wiki/Novikov_conjecture}), Signature (topologie) --- Wikipédia (en anglais) (\url{https://en.wikipedia.org/wiki/Signature_(topology)})
\item Article : A. Connes et H. Moscovici, \og{} Cyclic cohomology, the Novikov conjecture and hyperbolic groups \fg{} (1990), \textit{Topology} 29
\item Article : G. Yu, \og{} The coarse Baum--Connes conjecture for spaces which admit a uniform embedding into Hilbert space \fg{} (2000), \textit{Inventiones Mathematicae} 139
\end{refs}

\bigskip

\begin{problem}[{Le dernier invariant de Kervaire --- Recherche, short}]
\textbf{TOP-45.} (Résolu très récemment --- voir la note.) Une variété lisse close
parallélisée de dimension $ n = 4k+2 $ porte un invariant de Kervaire dans $ \mathbb{Z}/2 $, construit
comme l'invariant d'Arf d'un raffinement quadratique de la forme d'intersection en dimension médiane.
Déterminez tous les $ n $ pour lesquels la valeur $ 1 $ se produit effectivement, et rédigez un exposé
honnête de la façon dont la réponse a été assemblée et de ce qui repose encore sur des travaux très
récents.
\end{problem}

\noindent Kervaire s'est servi de l'invariant en 1960 pour construire une variété de dimension
$ 10 $ ne portant aucune structure lisse. Browder a démontré en 1969 que la valeur $ 1 $ ne peut se
produire qu'en dimensions $ 2^{j} - 2 $, et des constructions étaient connues en dimensions
$ 2, 6, 14, 30 $ puis --- après les travaux de Barratt, Jones et Mahowald dans les années 1980 ---
$ 62 $. Hill, Hopkins et Ravenel ont ensuite démontré (annonce en 2009 ;
\textit{Annals of Mathematics} 184, 2016) qu'elle ne se produit jamais en dimension $ 254 $ ou au-delà,
laissant exactement un cas non tranché : $ 126 $. En décembre 2024 Lin, Wang et Xu ont déposé une
démonstration que la dimension $ 126 $ contient bien une classe d'invariant de Kervaire un, complétant la
liste --- un calcul de la suite spectrale d'Adams mené avec une assistance machine substantielle. En 2026 cet
argument est accepté par les spécialistes et a passé la relecture par les pairs, mais il est jeune ;
énoncez-le comme réglé et dites qui l'a réglé.

\begin{refs}
\item Contexte : Invariant de Kervaire --- Wikipédia (en anglais) (\url{https://en.wikipedia.org/wiki/Kervaire_invariant}), Suite spectrale d'Adams --- Wikipédia (en anglais) (\url{https://en.wikipedia.org/wiki/Adams_spectral_sequence})
\item Article : M. A. Hill, M. J. Hopkins, D. C. Ravenel, \og{} On the non-existence of elements of Kervaire invariant one \fg{} (2016), \textit{Annals of Mathematics} 184, arXiv:0908.3724 (\url{https://arxiv.org/abs/0908.3724})
\item Article : W. Lin, G. Wang, Z. Xu, \og{} On the Last Kervaire Invariant Problem \fg{} (2024), arXiv:2412.10879 (\url{https://arxiv.org/abs/2412.10879})
\end{refs}

\bigskip

\begin{problem}[{La conjecture de Borel --- Recherche}]
\textbf{TOP-46.} (Ouvert.) Une variété close est \textit{asphérique} si son revêtement universel est
contractile --- autrement dit si $ \pi_n = 0 $ pour tout $ n \ge 2 $. La conjecture de Borel affirme
qu'une telle variété est déterminée à homéomorphisme près par son seul groupe fondamental : toute
équivalence d'homotopie entre variétés closes asphériques est homotope à un homéomorphisme.
\begin{enumerate}
\item[(a)] Démontrez la moitié facile de la correspondance : deux variétés closes asphériques de groupes
fondamentaux isomorphes sont homotopiquement équivalentes. (Les applications vers un CW-complexe
asphérique sont déterminées à homotopie près par ce qu'elles font à $ \pi_1 $.)
\item[(b)] Réglez la dimension deux à la main : à l'aide de la classification des surfaces (TOP-19) et des
groupes de surfaces abélianisés de TOP-18(c), démontrez que deux surfaces closes asphériques de groupes
fondamentaux isomorphes sont homéomorphes. Lisez ensuite pourquoi l'énoncé de rigidité plus fort vaut
également pour les surfaces, par le théorème de Dehn--Nielsen--Baer.
\item[(c)] Montrez que l'analogue lisse est faux : si $ \Sigma $ est une sphère exotique, la somme connexe
$ M \# \Sigma $ est homéomorphe à $ M $ mais n'a pas besoin de lui être difféomorphe. Expliquez
pourquoi cela fait d'\og{} homéomorphisme \fg{}, et non de \og{} difféomorphisme \fg{}, le bon mot dans la
conjecture.
\item[(d)] La conjecture elle-même : démontrez que toute équivalence d'homotopie de variétés closes
asphériques est homotope à un homéomorphisme, ou produisez un contre-exemple. Comme voie d'entrée,
rédigez un panorama du programme de Farrell--Jones et de la question de savoir exactement quelles classes
de groupes fondamentaux il couvre.
\end{enumerate}

\end{problem}

\noindent Armand Borel a formulé la conjecture vers 1953, dans sa correspondance avec Serre,
après avoir vu la rigidité à la Mostow pour les espaces localement symétriques et s'être demandé si la
géométrie était vraiment nécessaire. On en sait beaucoup : les dimensions au plus trois découlent de la
géométrisation (Perelman, 2003) ; Farrell et Jones l'ont démontrée en 1989--93 pour les variétés
riemanniennes closes à courbure négative ou nulle de dimension au moins cinq ; et Bartels et Lück ont
démontré la conjecture de Farrell--Jones en $ K $- et $ L $-théorie pour les groupes hyperboliques au
sens des mots et pour les groupes CAT(0) (\textit{Annals of Mathematics} 175, 2012), ce qui donne Borel en
dimension au moins cinq pour tous ces groupes fondamentaux. La même machinerie d'applications d'assemblage
livre aussi la conjecture de Novikov (TOP-36), et c'est pourquoi les deux problèmes font route ensemble. En
2026 la conjecture générale est ouverte et aucun contre-exemple n'est connu.

\begin{refs}
\item Contexte : Conjecture de Borel --- Wikipédia (en anglais) (\url{https://en.wikipedia.org/wiki/Borel_conjecture}), Espace asphérique --- Wikipédia (\url{https://fr.wikipedia.org/wiki/Espace_asph%C3%A9rique})
\item Contexte : Conjecture de Farrell--Jones --- Wikipédia (en anglais) (\url{https://en.wikipedia.org/wiki/Farrell%E2%80%93Jones_conjecture})
\item Article : A. Bartels et W. Lück, \og{} The Borel Conjecture for hyperbolic and CAT(0)-groups \fg{} (2012), \textit{Annals of Mathematics} 175, arXiv:0901.0442 (\url{https://arxiv.org/abs/0901.0442})
\item Lecture : A. Ranicki, \textit{Algebraic and Geometric Surgery}, Oxford University Press
\end{refs}

\bigskip


\vfill
\noindent\emph{Hors de la Caverne} --- des probl\`emes, pas de r\'eponses.
\end{document}
